\subsection{连续有限维表示}

设 \(\varrho\) 是拓扑群 G 在有限维分离拓扑 K-向量空间 E 中的连续线性表示。我们赋予 \(\mathcal{L}(\mathrm{E})\) 其在 K 上唯一的分离拓扑向量空间结构；于是，态射 \(\varrho\colon\mathrm{G}\to\mathrm{GL}(\mathrm{E})\) 连续，因为 \(\mathrm{GL}(\mathrm{E})\) 的拓扑由 \(\mathcal{L}(\mathrm{E})\) 的拓扑诱导，而该拓扑与逐点收敛拓扑重合（EVT, I, p. 14, th. 2）。因此，表示 \(\varrho\) 是连续的。如果 G 是实李群，那么 \(\varrho\) 的特征标是 G 上的解析函数（LIE, III, p. 225, § 8, no. 1, th. 1）。当 E' 赋予其作为 K 上分离拓扑向量空间的唯一拓扑时，逆步表示也是连续的。此外，对每个整数 \(n\geqslant0\)，当相应空间赋予其作为 K 上分离拓扑向量空间的拓扑时，表示 \(\mathrm{T}^{n}(\varrho)\) 、\(\mathrm{S}^{n}(\varrho)\) 和 \(\wedge^{n}(\varrho)\) （同上）都是连续的。

设 \(\varrho_{1}\) 和 \(\varrho_{2}\) 是拓扑群 G 在有限维分离拓扑向量空间 E \(_{1}\) 和 E \(_{2}\) 中的连续表示。由 \(\varrho_{1} \otimes \varrho_{2}\) 给出的 G 在 E \(_{1} \otimes\) E \(_{2}\) 中的线性表示（LIE, III, p. 256, 附录）是连续的；空间 E \(_{1} \otimes\) E \(_{2}\) 赋予其作为 K 上分离拓扑向量空间的拓扑。

\subsection{不可约表示}

本节中，G 是拓扑群。

定义 2. --- 拓扑 K-向量空间 E 中 G 的表示 \(\varrho\) 称为不可约的，是指 \(\{0\}\) 在 E 中不稠密，并且 \(\varrho\) 的唯一子表示是 \(\varrho\) 本身和在 \(\{0\}\) 的闭包中诱导的表示。

如果 \(\varrho\) 是分离拓扑向量空间 E 中 G 的不可约表示，那么 E 的每个非零元素都是 \(\varrho\) 的循环向量。

引理 2. --- 设 \(\pi\) 和 \(\varrho\) 是分离拓扑 K-向量空间中 G 的连续线性表示。假设 \(\pi\) 不可约。从 \(\pi\) 到 \(\varrho\) 的每个非零 G-态射都是单射，而从 \(\varrho\) 到 \(\pi\) 的每个非零 G-态射都有稠密像。

特别地，如果 \(\pi\) 和 \(\varrho\) 都不可约且有限维，那么从 \(\pi\) 到 \(\varrho\) 的每个非零 G-态射都是同构。

由于 \(\varrho\) 的表示空间是分离的，G-态射 \(u\)（从 \(\pi\) 到 \(\varrho\)）的核是闭的，因而定义 \(\pi\) 的一个子表示。如果态射 \(u\) 非零，其核必为 \(\{0\}\)，因为 \(\pi\) 是分离拓扑向量空间中的不可约表示。同样，从 \(\varrho\) 到 \(\pi\) 的非零 G-态射的像的闭包是 \(\pi\) 的一个非零子表示，因而等于 \(\pi\) 的表示空间。

最后一个断言由上述内容推出。

引理 3. --- 设 \(\varrho\) 是分离拓扑 K-向量空间 E 中的有限维连续线性表示。如果 E 非零，那么 \(\varrho\) 存在一个不可约子表示。

由于 E 有限维且非零，存在维数最小的非零 G-不变子空间 F。该子空间在 E 中闭，并定义 E 的一个子表示；F 的每个子表示也是 E 的子表示，因此由极小性，F 中的表示不可约。

注. --- G 的一个非零表示 E 并不总是包含不可约子表示（参见 V, p. 426, 注）。不过将会看到，当 G 是紧群、K = C，且 E 是 K 上非零的分离拟完备局部凸空间时，情况确实如此（V, p. 464 的命题 7）。

\subsection{酉表示}

定义 3. --- 设 G 是拓扑群，E 是 K 上的希尔伯特空间。G 在 E 中的酉表示，是 G 在 E 中的连续线性表示 \(\varrho\)，使得对 G 中每个 g，E 的自同态 \(\varrho(g)\) 都是 E 的酉自同态（EVT, V, p. 40）。

换言之，酉表示是希尔伯特空间中的等距表示。特别地，酉表示是有界的。

拓扑群在希尔伯特空间中的平凡表示是酉的。酉表示的每个子表示都是酉的。酉表示对子群的限制是酉的。

定义 4. --- 设 G 是拓扑群，\(\varrho\) 和 \(\pi\) 分别是 G 在希尔伯特空间 E 和 F 中的酉表示。E × F 上的 G-不变半双线性形式，是 E × F 上的连续半双线性形式 q，使得

\[
q (x, y) = q (\varrho (g) x, \pi (g) y)
\]

对所有 \(g \in \mathrm{G}\) 及每个 \((x, y) \in \mathrm{E} \times \mathrm{F}\) 成立。

E × F 上的 G-不变半双线性形式组成的向量空间记为 \(\mathrm{Sesq}_{\mathrm{G}}(\varrho,\pi)\) 或 \(\mathrm{Sesq}_{\mathrm{G}}(\mathrm{E},\mathrm{F})\)。

例. --- 设 G 是拓扑群，\(\varrho\) 是 G 在 E 中的酉表示。E 上的内积 \(q\) 是 \(\mathrm{E} \times \mathrm{E}\) 上的 G-不变半双线性形式。

引理 4. --- 设 G 是拓扑群，\(\varrho\) 是 G 到希尔伯特空间 E 的酉群 U(E) 中的同态。那么，\(\varrho\) 是 G 的酉表示，当且仅当对于 U(E) 上的逐点收敛拓扑，\(\varrho\) 在 G 的元素 e 处连续。只需该性质对 E 的一个总子集中的每个 \(x\) 成立即可。

该条件显然是必要的。根据 INT, VIII, p. 129, § 2, no. 1 的注 2，它也是充分的，因为 \(\varrho\) 的像在 \(\mathcal{L}(\mathrm{E})\) 中等度连续，并且映射 \(g \mapsto \varrho(g)x\) 在 \(e\) 处的连续性蕴含其在 G 上的连续性。

设 G 是拓扑群。如果 \(\varrho_{1}\) 和 \(\varrho_{2}\) 是 G 的酉表示，且 u 属于 \(\operatorname{Hom}_{\mathrm{G}}(\varrho_{1},\varrho_{2})\)，那么 \(u^{*}\) 属于 \(\operatorname{Hom}_{\mathrm{G}}(\varrho_{2},\varrho_{1})\)。事实上，由于 \(\varrho_{1}\) 和 \(\varrho_{2}\) 是酉表示，对每个 \(g\in G\) 有

\[
\begin{array}{r} u ^ {*} \circ \varrho_ {2} (g) = u ^ {*} \circ \varrho_ {2} (g ^ {- 1}) ^ {*} = (\varrho_ {2} (g ^ {- 1}) \circ u) ^ {*} \\ = (u \circ \varrho_ {1} (g ^ {- 1})) ^ {*} = \varrho_ {1} (g) \circ u ^ {*}. \end{array}
\]

引理 5. --- 设 G 是拓扑群，\(\varrho\) 是 G 在复希尔伯特空间 E 中的酉表示。空间 \(\mathrm{Hom}_{\mathrm{G}}(\varrho, \varrho)\) 是 \(\mathcal{L}(\mathrm{E})\) 的含单位星子代数。

前述内容表明，\(\operatorname{Hom}_{\mathrm{G}}(\varrho,\varrho)\) 是 \(\mathcal{L}(\mathrm{E})\) 的含单位自伴子代数。由于它在 \(\mathcal{L}(\mathrm{E})\) 中是闭的，所以它是 \(\mathcal{L}(\mathrm{E})\) 的星子代数。

引理 6. --- 设 \(\pi\) 和 \(\varrho\) 分别是拓扑群 G 在希尔伯特空间 E 和 F 中的酉表示。设 D 是 E 的稠密子空间，并且在 \(\pi\) 下稳定。设 u 是从 E 到 F、定义域为 D 的闭部分算子，并满足关系式 \(u \circ \pi(g) = \varrho(g) \circ u\) 对所有 \(g \in G\) 成立。则 \(u^{*}\) 的定义域在 \(\varrho\) 下稳定，并且有关系式 \(u^{*} \circ \varrho(g) = \pi(g) \circ u^{*}\) 对所有 \(g \in G\) 成立。

设 \(g \in \mathrm{G}\)，且 \(x \in \mathrm{dom}(u^*)\)。对于每个 \(y \in \mathrm{dom}(u)\)，有

\[
\begin{array}{r l} & {\langle \varrho (g) x \mid u (y) \rangle = \langle x \mid \varrho (g ^ {- 1}) u (y) \rangle = \langle x \mid u (\pi (g ^ {- 1}) y) \rangle} \\ & {\qquad = \langle u ^ {*} (x) \mid \pi (g ^ {- 1}) y \rangle = \langle \pi (g) (u ^ {*} (x)) \mid y \rangle ,} \end{array}
\]

因为 \(\varrho(g)^{*} = \varrho(g)^{-1} = \varrho(g^{-1})\)。这证明了 \(\varrho(g)x \in \mathrm{dom}(u^{*})\)，且 \(u^{*}(\varrho(g)x) = \pi(g)(u^{*}(x))\)。特别地，\(u^{*} \circ \varrho(g)\) 的定义域包含 \(u^{*}\) 的定义域，并且 \(\pi(g) \circ u^{*} \subset u^{*} \circ \varrho(g)\)。

此外，若 \(x \in \text{dom}(u^{*} \circ \varrho(g))\)，则 \(x = \varrho(g^{-1})(\varrho(g)x)\) 属于 \(\text{dom}(u^{*})\)，这是将前述结论应用于 \(g^{-1}\) 的结果。因此 \(u^{*} \circ \varrho(g) = \pi(g) \circ u^{*}\)。

命题 1. --- 设 \(\varrho_{1}\) 和 \(\varrho_{2}\) 分别是拓扑群 G 在希尔伯特空间 E \(_{1}\) 和 E \(_{2}\) 中的酉表示。从 Hom \(_{G}(\varrho_{1},\varrho_{2})\) 到 Sesq \(_{G}(\varrho_{2},\varrho_{1})\) 的映射将 u 关联到半双线性形式 \(q_{u}\)，该形式由 \(q_{u}(x,y)=\langle x|u(y)\rangle\) 定义；它是向量空间的同构。

若 \(u\) 是 G-态射，则有

\[
\begin{array}{c} q _ {u} (\varrho_ {2} (g) x, \varrho_ {1} (g) y) = \langle \varrho_ {2} (g) x \mid u (\varrho_ {1} (g)y) \rangle \\ = \langle \varrho_ {2} (g) x \mid \varrho_ {2} (g) u (y) \rangle = \langle x \mid u (y) \rangle = q _ {u} (x, y) \end{array}
\]

对所有 \(g \in G\) 及每个 \((x, y) \in \mathrm{E}_{2} \times \mathrm{E}_{1}\) 成立，因此所指映射是从 \(\operatorname{Hom}_{\mathrm{G}}(\varrho_{1}, \varrho_{2})\) 到 \(\operatorname{Sesq}_{\mathrm{G}}(\varrho_{2}, \varrho_{1})\) 的线性映射。由 EVT, V, p. 16, cor. 2，它是单射。

反之，设 \(q\) 是 G-不变半双线性形式，且 \(u\) 是从 \(\mathrm{E}_1\) 到 \(\mathrm{E}_2\) 的唯一线性映射，使得 \(q(x,y) = \langle x \mid u(y) \rangle\) 对所有 \((x,y) \in \mathrm{E}_2 \times \mathrm{E}_1\) 成立（同上）。对于每个 \(g\)（\(G\) 中）及每个 \((x,y)\)

在 \(\mathrm{E}_2\times \mathrm{E}_1\) 中，我们有

\[
\begin{array}{c} \langle x | (\varrho_ {2} (g) \circ u \circ \varrho_ {1} (g ^ {- 1})) (y) \rangle = \langle \varrho_ {2} (g ^ {- 1}) x | u (\varrho_ {1} (g ^ {- 1}) y) \rangle \\ = q (\varrho_ {2} (g ^ {- 1}) x, \varrho_ {1} (g ^ {- 1}) y) = q (x, y) = \langle x | u (y) \rangle , \end{array}
\]

因此 \(u = \varrho_{2}(g) \circ u \circ \varrho_{1}(g^{-1})\)。于是 \(u \in \operatorname{Hom}_{\mathrm{G}}(\varrho_{1}, \varrho_{2})\)。命题得证。

设 \(\varrho\) 是拓扑群 G 在希尔伯特空间 E 中的酉表示。若 K = C，记 \(\overline{\mathrm{E}}\) 为 E 的共轭空间（EVT, V, p. 6）。若 K = R，令 \(\overline{\mathrm{E}} = \mathrm{E}\)。共轭表示 \(\overline{\varrho}\) 是 G 在 \(\overline{\mathrm{E}}\) 中的表示，由 \(\overline{\varrho}(g) = \varrho(g)\)（对所有 \(g \in \mathrm{G}\)）定义。它是 G 的酉表示；按定义，\(\overline{\mathrm{E}}\) 的一个子空间是 \(\overline{\mathrm{E}}\) 的子表示，当且仅当它是 E 的子表示。

命题 2. --- 设 \(\varrho\) 是拓扑群 G 在希尔伯特空间 E 中的酉表示。记 u 为从 \(\overline{\mathrm{E}}\) 到 \(\mathrm{E}'\) 的等距同构，它将 \(x\) 关联到线性形式 \(y\mapsto \langle x|y\rangle\)。

\begin{enumerate}
\def\labelenumi{\alph{enumi})}
\item
  映射 u 是从共轭表示 \(\overline{\varrho}\) 到逆步表示 \(\breve{\varrho}\) 的同构；
\item
  给 E' 赋予通过 u 传递结构得到的希尔伯特空间结构；逆步表示 \(\breve{\varrho}\) 是 G 在 E' 中的酉表示。
\end{enumerate}

由 EVT, V, p. 15, th. 3 及其后注记，映射 \(u\) 是等距同构。

对 G 中每个 g、\(\overline{E}\) 中每个 x 以及 E 中每个 y，有

\[
\begin{array}{c} \langle y, (\breve {\varrho} (g) \circ u) (x) \rangle = \langle \varrho (g ^ {- 1}) y, u (x) \rangle \\ = \langle x | \varrho (g ^ {- 1}) y \rangle = \langle \varrho (g) x | y \rangle = \langle y, u (\overline {{\varrho}} (g) x) \rangle , \end{array}
\]

因此 \(\breve{\varrho}(g) \circ u = u \circ \overline{\varrho}(g)\)，这证明了 \(a\)；断言 \(b\) 立即得到。

推论. --- 设 \(\varrho\) 是拓扑群 G 在希尔伯特空间 E 中的酉表示。以下条件等价：

\begin{enumerate}
\def\labelenumi{(\roman{enumi})}
\item
  表示 \(\varrho\) 不可约；
\item
  逆步表示 \(\breve{\varrho}\) 不可约；
\item
  共轭表示 \(\overline{\varrho}\) 不可约。
\end{enumerate}

这由命题 2 以及命题前关于 \(\overline{E}\) 的子表示的注记推出。

命题 3. --- 设 \(\varrho\) 是拓扑群 G 在希尔伯特空间 E 中的酉表示。设 \(\pi\) 是 \(\varrho\) 的子表示，其表示空间为 F。F 在 E 中的正交补 F° 是 E 的子表示，且 E = F ⊕ F°，F° 与 E/F 同构。

空间 F° 是闭的。对于每个 \(x \in F^{\circ}\)、所有 \(g \in G\) 以及每个 \(y \in F\)，有 \(\langle \varrho(g)x \mid y \rangle = \langle x \mid \varrho(g^{-1})y \rangle = 0\)，因为 \(\varrho\) 是酉表示且 F 是子表示。因此 \(\varrho(g)x \in F^{\circ}\)。

E 到 E/F 的规范投影是 G-态射，所以它在 F° 上的限制是从 F° 到 E/F 的 G-态射；由 EVT, V, p. 13，它是等距同构。

我们也把 G 在 F° 中的表示记作 \(\pi^{\circ}\)。

命题 4. --- 设 \(\varrho\) 是拓扑群 G 在希尔伯特空间 E 中的酉表示。E 的闭子空间 F 是 \(\varrho\) 的子表示，当且仅当 E 上像为 F 的正交投影 p 是从 \(\varrho\) 到 \(\varrho\) 的 G-态射。

设 F 是 \(\varrho\) 的子表示。取 \(x \in E\)，写成 \(x = p(x) + y\)，其中 \(y \in F^{\circ}\)。对于 G 中每个 \(g\)，有

\[
\varrho (g) x = \varrho (g) (p (x)) + \varrho (g) y,
\]

并且由于 \(\varrho(g)(p(x))\) 属于 F 而 \(\varrho(g)y\) 属于 F°（命题 3），有 \(p(\varrho(g)x) = \varrho(g)(p(x))\)。因此 \(p\) 属于 \(\operatorname{Hom}_{\mathrm{G}}(\varrho, \varrho)\)。

反之，若 \(p \in \operatorname{Hom}_{\mathrm{G}}(\varrho, \varrho)\)，则 \(1_{E} - p\) 是 G-态射，因而 \(\mathrm{F} = \operatorname{Ker}(1_{\mathrm{E}} - p)\) 是 \(\varrho\) 的子表示。

\subsection{酉表示的希尔伯特直和与张量积}

设 G 是拓扑群。设 \((\varrho_{i})_{i\in\mathrm{I}}\) 是 G 在希尔伯特空间 \(E_{i}\) 中的一族酉表示。设 E 是各 \(E_{i}\) 的外部希尔伯特直和空间。对于 G 中每个 g 以及 E 中每个 \(x=(x_{i})_{i\in\mathrm{I}}\)，有

\[
\sum_ {i} \| \varrho_ {i} (g) x _ {i} \| ^ {2} = \sum_ {i} \| x _ {i} \| ^ {2} = \| x \| ^ {2},
\]

这证明了 \((\varrho_i(g)x_i)_{i\in \mathrm{I}}\) 属于 E 且与 \(x\) 有相同的范数。因此，元素 \(\varrho (g)\colon (x_{i})_{i\in \mathrm{I}}\mapsto (\varrho_{i}(g)x_{i})_{i\in \mathrm{I}}\) 是 \(\mathcal{L}(\mathrm{E})\) 中的酉元素。

引理 7. --- 映射 \(g \mapsto \varrho(g)\) 是 G 在 E 中的酉表示。

由 V, p. 380 的引理 4，只需证明对于 I 中每个 i 及 \(E_{i}\) 中每个 x，映射 \(g \mapsto \varrho_{i}(g)x\) 在 G 的单位元 e 处连续。该映射是连续映射 \(g \mapsto \varrho_{i}(g)x\) 与 \(E_{i}\) 到 E 的规范嵌入的复合，因而连续。

表示 \(\varrho\) 称为酉表示 \((\varrho_{i})_{i\in I}\) 的希尔伯特直和，记为 \(\varrho = \bigoplus \varrho_{i}\)。

设 G 和 H 是拓扑群。设 \(\varrho_{1}\)（分别为 \(\varrho_{2}\)）是 G（分别为 H）在希尔伯特空间 E \(_{1}\)（分别为 E \(_{2}\)）中的酉表示。设 E = E \(_{1}\) \(\widehat{\otimes}_{2}\) E \(_{2}\) 是 E \(_{1}\) 与 E \(_{2}\) 的希尔伯特张量积空间（EVT, V, p. 28, 定义 1）。对于 \((g,h)\in\mathrm{G}\times\mathrm{H}\)，令 \(\varrho(g,h)\) 为 E 的连续自同态 \(\varrho_{1}(g)\widehat{\otimes}_{2}\varrho_{2}(h)\)（TVS, V, p. 28）；无歧义时仍简记为 \(\varrho_{1}(g)\otimes\varrho_{2}(h)\)。

引理 8. --- a) 映射 \(\varrho\colon(g,h)\mapsto\varrho(g,h)\) 是 G × H 在 E 中的酉表示；

\begin{enumerate}
\def\labelenumi{\alph{enumi})}
\setcounter{enumi}{1}
\item
  对于任意标准正交基 \((e_i)_{i\in \mathrm{I}}\)，它是 \(\mathrm{E}_1\) 的基，从希尔伯特直和 \(\bigoplus_{i\in \mathrm{I}}\mathrm{E}_2\) 到 E 的映射 \((y_i)_{i\in \mathrm{I}}\mapsto \sum_{i\in \mathrm{I}}e_i\otimes y_i\) 是希尔伯特直和 \(\bigoplus_{i\in \mathrm{I}}\varrho_2\) 在 \(\operatorname {Res}_{\{e\} \times \mathrm{H}}^{\mathrm{G}\times \mathrm{H}}(\varrho)\) 中的等距 H-同构；
\item
  对于任意标准正交基 \((f_j)_{j\in \mathrm{J}}\)，它是 \(\mathrm{E}_2\) 的基，从希尔伯特直和 \(\bigoplus_{j\in \mathrm{J}}\mathrm{E}_1\) 到 E 的映射 \((x_{j})_{j\in \mathrm{J}}\mapsto \sum_{j\in \mathrm{J}}x_{j}\otimes f_{j}\) 是希尔伯特直和 \(\bigoplus_{j\in \mathrm{J}}\varrho_1\) 在 \(\operatorname {Res}_{\mathbf{G}\times \{e\}}^{\mathbf{G}\times \mathbf{H}}(\varrho)\) 中的等距 G-同构。
\end{enumerate}

映射 \((g,h)\mapsto\varrho(g,h)\) 是从 \(\mathrm{G}\times\mathrm{H}\) 到 \(\mathbf{GL}(\mathrm{E})\) 的同态（参见 EVT, V, p. 28, no. 2）。设 \((e_{i})_{i\in\mathrm{I}}\) 是 \(\mathrm{E}_{1}\) 的标准正交基。由 EVT, V, p. 29, prop. 3 和 cor. 2，映射

\[
u \colon (y _ {i}) \mapsto \sum_ {i \in I} e _ {i} \otimes y _ {i}
\]

是从希尔伯特直和 \(F_{2} = \bigoplus_{i \in I} E_{2}\) 到 E 上的等距映射。借助此等距映射，H 在 E 中的表示 \(\varrho_{\mathrm{H}} \colon h \mapsto \varrho(e, h)\) 被认同为表示族 \((\varrho_{2})_{i\in I}\) 在 \(F_{2}\) 中的直和表示。特别地，它是 H 在 E 中的酉表示（引理 7）。同样，同态 \(\varrho_{G}: g \mapsto \varrho(g,e)\) 是 G 在 E 中的酉表示。

取 \((g,h)\in\mathrm{G}\times\mathrm{H}\)。有 \(\varrho(g,h)=\varrho_{\mathrm{G}}(g)\circ\varrho_{\mathrm{H}}(h)\)，故 \(\varrho(g,h)\) 是酉的；此外，\(\varrho_{\mathrm{G}}(g)\) 与 \(\varrho_{\mathrm{H}}(h)\) 交换，所以由 V, p. 377 的引理 1，\(\varrho\) 是 \(G\times H\) 的酉表示。

最后，关于 \(\varrho\) 限制到子群 \(\{e\} \times \mathrm{H}\) 和 \(\mathrm{G} \times \{e\}\) 的断言已在上述论证过程中得到。

G×H 的表示 \((g,h)\mapsto\varrho_{1}(g)\otimes\varrho_{2}(h)\) 称为酉表示 \(\varrho_{1}\) 与 \(\varrho_{2}\) 的外张量积，记为 \(\varrho_{1}\boxtimes\varrho_{2}\)。

设 \(n \in \mathbf{N}\)，且 \((\mathrm{G}_i)_{1 \leqslant i \leqslant n}\) 是一族有限个拓扑群。设 \(\varrho_i\) 是 \(\mathrm{G}_i\) 在希尔伯特空间 \(\mathrm{E}_i\) 中的酉表示，对 \(1 \leqslant i \leqslant n\) 成立。类似地定义表示

\[
\varrho_ {1} \boxtimes \dots \boxtimes \varrho_ {n}
\]

\(G_{1} \times \cdots \times G_{n}\) 在希尔伯特空间 \(E = E_{1} \widehat{\otimes}_{2} \cdots \widehat{\otimes}_{2} E_{n}\) 中的表示（EVT, V, p. 27）。

若 \(G_{i} = G\)（\(1 \leqslant i \leqslant n\)），记 \(\Delta_{n}: G \to G^{n}\) 为由 \(g \mapsto (g, \ldots, g)\)（对所有 \(g \in G\)）定义的同态。我们用 \(\varrho_{1} \otimes \cdots \otimes \varrho_{n}\) 表示 G 的酉表示 \((\varrho_{1} \boxtimes \cdots \boxtimes \varrho_{n}) \circ \Delta_{n}\)。这称为酉表示 \(\varrho_{i}\) 的张量积。

对于每个置换 \(\sigma\)，即 \(\{1,\ldots,n\}\) 上的置换，规范同构

\[
\mathrm{E} _ {1} \widehat {\otimes} _ {2} \dots \widehat {\otimes} _ {2} \mathrm{E} _ {n} \rightarrow \mathrm{E} _ {\sigma (1)} \widehat {\otimes} _ {2} \dots \widehat {\otimes} _ {2} \mathrm{E} _ {\sigma (n)}
\]

（EVT, V, p. 28）是等距同构

\[
\varrho_ {1} \otimes \dots \otimes \varrho_ {n} \rightarrow \varrho_ {\sigma (1)} \otimes \dots \otimes \varrho_ {\sigma (n)}
\]

且是 G 的表示之间的同构。

若 \(\varrho_{i}=\varrho\) 对 \(1\leqslant i\leqslant n\) 成立，其中 \(\varrho\) 是 G 的酉表示，我们也用 \(\varrho^{\otimes n}\) 表示各 \(\varrho_{i}\) 的张量积表示，并称其为 \(n\) 次张量幂，即 \(\varrho\) 的张量幂。

\subsection{矩阵系数}

定义 5. --- 设 G 是拓扑群，\(\varrho\) 是 G 在希尔伯特空间 E 中的酉表示。设 \(x\) 和 \(y\) 是 E 中的元素。由 \(g \mapsto \langle x | \varrho(g)y \rangle\) 给出的从 G 到 K 的函数称为 \(\varrho\) 的矩阵系数，也称代表函数。若 \(x = y\)，则称其为对角矩阵系数。若 \(\varrho\) 是有限维的，则称其为有限维矩阵系数。

\(\varrho\) 的矩阵系数是 G 上的连续有界函数。我们用 \(\Upsilon(G)\)（分别用 \(\Theta(G)\)）表示 G 的复酉表示（分别为复有限维表示）的矩阵系数集合。

命题 5. — 设 G 是拓扑群。集合 \(\Theta(\mathrm{G})\) 和 \(\Upsilon(\mathrm{G})\) 是 \(\mathcal{C}_{\mathrm{b}}(\mathrm{G})\) 的含单位对合子代数。

常值函数 1 是 G 在 C 上的平凡表示的矩阵系数。设 \(\varrho\) 是 G 的酉表示，且 \(x\) 和 \(y\) 是 \(\varrho\) 的表示空间中的向量。对每个 \(\lambda \in \mathbf{C}\) 及每个 \(g \in \mathrm{G}\)，有 \(\lambda \langle x \mid \varrho(g)y \rangle = \langle x \mid \varrho(g)(\lambda y) \rangle\)。此外，有

\[
\langle x \mid \varrho (g) y \rangle = \langle \varrho (g ^ {- 1}) x \mid y \rangle = \overline {{\langle y \mid \varrho (g ^ {- 1}) x \rangle}}
\]

对所有 \(g \in G\) 成立。因此，集合 \(\Theta(G)\) 和 \(\Upsilon(G)\) 在标量乘法和共轭下稳定。

设 \(\varrho_{1}\) 和 \(\varrho_{2}\) 是 G 的酉表示；设 \((x_{1}, y_{1})\) 是 \(\varrho_{1}\) 的表示空间中的向量，\((x_{2}, y_{2})\) 是 \(\varrho_{2}\) 的表示空间中的向量。于是对每个 \(g \in G\)，有

\[
\begin{array}{c} \langle x _ {1}   |   \varrho_ {1} (g) y _ {1} \rangle + \langle x _ {2}   |   \varrho_ {2} (g) y _ {2} \rangle = \langle (x _ {1}, x _ {2})   |   (\varrho_ {1} \oplus \varrho_ {2}) (g) (y _ {1}, y _ {2}) \rangle , \\ \langle x _ {1}   |   \varrho_ {1} (g) y _ {1} \rangle \langle x _ {2}   |   \varrho_ {2} (g) y _ {2} \rangle = \langle x _ {1} \otimes x _ {2}   |   (\varrho_ {1} \otimes \varrho_ {2}) (g) (y _ {1} \otimes y _ {2}) \rangle , \end{array}
\]

这证明了 \(\Theta(G)\) 和 \(\Upsilon(G)\) 在加法和乘法下稳定。命题得证。

\subsection{舒尔引理}

本节中，希尔伯特空间均为复的。

命题 6（舒尔引理）。--- 设 \(\varrho\) 是拓扑群 G 在非零希尔伯特空间 E 中的酉表示。那么，\(\varrho\) 不可约，当且仅当 \(\mathrm{Hom}_{\mathrm{G}}(\varrho, \varrho)\) 等于 \(\mathbf{C} \cdot 1_{\mathrm{E}}\)。

\(\mathrm{Hom}_{\mathrm{G}}(\varrho, \varrho)\) 是 \(\mathcal{L}(\mathrm{E})\) 的含单位星子代数（V, p. 380 的引理 5）。对于 E 的每个子表示 F，像为 F 的正交投影 \(p\) 是星代数 \(\mathrm{Hom}_{\mathrm{G}}(\varrho, \varrho)\) 中的幂等元（V, p. 383 的命题 4）。若该代数等于 \(\mathbf{C} \cdot 1_{\mathrm{E}}\)，这意味着 \(p = 0\) 或 \(p = 1_{\mathrm{E}}\)，也即 F = 0 或 F = E。因此 \(\varrho\) 不可约。

反之，假设 \(\varrho\) 不可约。设 \(u\) 和 \(v\) 是星代数 \(\mathrm{Hom}_{\mathrm{G}}(\varrho, \varrho)\) 中满足 \(uv = 0\) 的交换元素。设 \(u\) 非零。于是 \(u\) 的核 F 定义 \(\varrho\) 的一个不同于 E 的子表示，故 F 必为 \(\{0\}\)；由于 F 包含 \(v\) 的像，推出 \(v = 0\)。因此由 I, p. 113 的命题 10，有 \(\mathrm{Hom}_{\mathrm{G}}(\varrho, \varrho) = \mathbf{C} \cdot 1_{\mathrm{E}}\)。

推论 1. --- 设 π 是拓扑群 G 在希尔伯特空间 E 中的不可约酉表示。设 u 是 E 上的闭部分算子。假设 u 的定义域稠密、dom(u) 在 π 下稳定，并且对所有 g ∈ G 有 u ∘ π(g) = π(g) ∘ u。那么 dom(u) = E，且 u 是恒等映射的标量倍数。

部分算子 \(u^{*} \circ u\) 是 E 上的正自伴部分算子（IV, p. 241 的命题 12）；对于 \(v = 1_{\mathrm{E}} + u^{*} \circ u\) 也同样如此，并且后者是单射的，因为 \(-1 \notin \operatorname{Sp}(u^{*} \circ u)\)（IV, p. 248 的命题 17）。有关系式 \(u^{*} \circ \pi(g) = \pi(g) \circ u^{*}\) 对所有 \(g \in G\) 成立（V, p. 381 的引理 6）；从而关系式 \(v \circ \pi(g) = \pi(g) \circ v\) 对所有 \(g \in G\) 成立。由于 \(v\) 是单射，故关系式 \(v^{-1} \circ \pi(g) = \pi(g) \circ v^{-1}\) 对所有 \(g \in G\) 成立。由于 \(v^{-1}\) 属于 \(\mathcal{L}(\mathrm{E})\)，所以由命题 6，存在 \(\lambda \in \mathbf{C}\) 使得 \(v^{-1} = \lambda 1_{\mathrm{E}}\)。必有 \(\lambda \neq 0\)，这意味着 \(\mathrm{E} = \mathrm{Im}(v^{-1}) = \mathrm{dom}(v) \subset \mathrm{dom}(u)\)，因而 \(\mathrm{dom}(u) = \mathrm{E}\)。由于 \(u\) 是闭的，有 \(u \in \mathcal{L}(\mathrm{E})\)，于是 \(u \in \mathrm{Hom}_{\mathrm{G}}(\pi, \pi)\)，并且由命题 6，\(u\) 是一个伸缩变换。

推论 2. --- 设 \(\varrho\) 和 \(\pi\) 是拓扑群 G 在希尔伯特空间 E 和 F 中的不可约酉表示。空间 \(\mathrm{Hom}_{\mathrm{G}}(\varrho, \pi)\) 的维数在 \(\varrho\) 与 \(\pi\) 同构时为 1，否则为零。特别地，若 \(\varrho\) 与 \(\pi\) 同构，则从 \(\varrho\) 到 \(\pi\) 的每个非零 G-态射都是同构。

假设存在非零 G-态射 \(u\)，它从 \(\varrho\) 到 \(\pi\)。线性映射 \(u^{*} \circ u\) 是 \(\mathrm{Hom}_{\mathrm{G}}(\varrho, \varrho)\) 的元素；因此存在复数 \(\lambda\)，使得 \(u^{*} \circ u = \lambda \cdot 1_{\mathrm{E}}\)（命题 6）。

于是有

\[
\langle u (x) \mid u (y) \rangle = \langle x \mid u ^ {*} u (y) \rangle = \lambda \langle x \mid y \rangle
\]

对 E 中所有 \(x\) 和 \(y\) 成立。特别地，\(\lambda \neq 0\)，因为 \(u\) 非零。由于 \(\| u(x)\| = |\lambda |^{1 / 2}\| x\|\) 对所有 \(x\in \mathrm{E}\) 成立，线性映射 \(u\) 是单射，且 \(u\) 的像在 F 中闭（I, p. 107 的引理 8）；于是它是不可约表示 \(\pi\) 的一个非零子表示，从而 \(u\) 是满射。因此，\(u\) 是从 \(\varrho\) 到 \(\pi\) 的同构。于是映射 \(v\mapsto u^{*}\circ v\) 是从空间 \(\mathrm{Hom}_{\mathrm{G}}(\varrho ,\pi)\) 到空间 \(\mathrm{Hom}_{\mathrm{G}}(\varrho ,\varrho)\) 的同构，而后一个空间的维数为 1（命题 6）。

推论 3. --- 设 \(\varrho\) 是拓扑群 G 在希尔伯特空间 E 中的酉表示。表示 \(\varrho\) 不可约，当且仅当空间 \(\mathrm{Sesq}_{\mathrm{G}}(\varrho, \varrho)\) 的维数为 1。此时，该空间由 E 上的内积生成。

根据 V, p. 381 的命题 1，这由命题 6 得出。

推论 4. --- 设 \(\varrho_{1}\) 和 \(\varrho_{2}\) 是拓扑群 G 在希尔伯特空间 E \(_{1}\) 和 E \(_{2}\) 中的非同构不可约酉表示。空间 Sesq \(_{G}\)（\(\varrho_{1}, \varrho_{2}\)）为零。

由 V, p. 381 的命题 1，这由推论 2 得出。

推论 5. --- 设 \(\varrho\) 是拓扑群 G 在希尔伯特空间 E 中的酉表示。设 \(\pi\) 是 G 在希尔伯特空间 F 中的不可约酉表示。

\begin{enumerate}
\def\labelenumi{\alph{enumi})}
\item
  从 \(\pi\) 到 \(\varrho\) 的每个非零 G-态射 u，都存在 \(\lambda \in \mathbf{R}_{+}^{*}\)，使得 \(\lambda u\) 是等距的；
\item
  从 π 到 ρ 的每个 G-态射 u 都有闭像；
\item
  从 \(\varrho\) 到 \(\pi\) 的每个非零 G-态射 v 都是满射。
\end{enumerate}

证明断言 \(a\)，它蕴含断言 \(b\)（I, p. 107 的引理 8）。由于态射 \(u\) 非零，它是单射（V, p. 378 的引理 2）。由公式 \(q(x,y) = \langle u(x)|u(y)\rangle\)，其中 \(x\) 和 \(y\) 属于 F，定义了 F 上的连续半双线性形式。由于 \(u\) 是 G-态射，有 \(q(\pi(g)x,\pi(g)y) = q(x,y)\) 对所有 \(g\)（属于 G）及 \((x,y)\in\mathrm{F}\times\mathrm{F}\) 成立，故 \(q\) 是 G-不变的。由推论 3，存在 \(\alpha\in\mathbf{R}_{+}^{*}\)，使得

\[
q (x, y) = \langle u (x) | u (y) \rangle = \alpha \langle x | y \rangle
\]

对所有 \((x,y)\in\mathrm{F}\times\mathrm{F}\) 成立，因而 \(\alpha^{-1/2}u\) 是等距的。

最后证明断言 \(c\)。由于 \(\pi\) 不可约，\(v\) 的像在 F 中稠密（V, p. 378 的引理 2），故伴随 \(v^*\) 是单射（TVS, V, p. 41 的命题 4）。由 \(b\)，\(v^*\) 的像 H 是 E 的闭子空间；因此它是 \(\varrho\) 的子表示，而由限制到子空间得到的 \(v^*\) 是从 F 到 H 的 G-同构。映射 \(w\) 是 \(v\) 在 H 上的限制，它定义了从 H 到 F 的 G-态射；由于其核是 H 与 \(\text{Ker}(v) = \text{H}^\circ\) 的交，它是单射（同上）。因此 \(w\) 是同构（推论 2）；特别地，\(v\) 是满射。

推论 6. --- 设 \(G_{1}\) 和 \(G_{2}\) 是拓扑群。设 \(\varrho_{1}\) 是 \(G_{1}\) 在希尔伯特空间 \(E_{1}\) 中的不可约酉表示，\(\varrho_{2}\) 是 \(G_{2}\) 在希尔伯特空间 \(E_{2}\) 中的不可约酉表示。外张量积 \(\varrho_{1} \boxtimes \varrho_{2}\) 是 \(G_{1} \times G_{2}\) 在 \(E_{1} \widehat{\otimes}_{2} E_{2}\) 中的不可约酉表示。

由 V, p. 384 的引理 8，外张量积 \(\varrho = \varrho_{1} \otimes \varrho_{2}\) 是 \(G = G_{1} \times G_{2}\) 在空间 \(E = E_{1} \widehat{\otimes}_{2} E_{2}\) 中的酉表示。

设 \(q\) 是 E 上的 \((\mathrm{G}_1 \times \mathrm{G}_2)\)-不变半双线性形式。对于每个 \((x_1, y_1) \in \mathrm{E}_1^2\)，映射 \((x_2, y_2) \mapsto q(x_1 \otimes x_2, y_1 \otimes y_2)\) 属于 \(\operatorname{Sesq}_{\mathrm{G}_2}(\varrho_2, \varrho_2)\)。记 \(b(x_1, y_1)\) 为满足下式的唯一复数：

\[
q (x _ {1} \otimes x _ {2}, y _ {1} \otimes y _ {2}) = b (x _ {1}, y _ {1}) \left\langle x _ {2} \mid y _ {2} \right\rangle
\]

对所有 \((x_{2},y_{2})\in\mathrm{E}_{2}^{2}\) 成立（推论 3）。取 \(\varepsilon\in E_{2}\) 使其范数为 1，于是有 \(b(x_{1},y_{1})=q(x_{1}\otimes\varepsilon,y_{1}\otimes\varepsilon)\) 对所有 \((x_{1},y_{1})\in\mathrm{E}_{1}^{2}\) 成立。该公式说明映射 b 是 \(E_{1}\) 上的半双线性映射。此外

\[
\| b (x _ {1}, y _ {1}) \| \leqslant \| q \| \| x _ {1} \otimes \varepsilon \| \| y _ {1} \otimes \varepsilon \| = \| q \| \| x _ {1} \| \| y _ {1} \|
\]

对所有 \((x_{1},y_{1})\in\mathrm{E}_{1}^{2}\) 成立（EVT, V, p. 26, 公式（5）），因此 b 是连续的。

设 \(g \in G\)，且 \((x_{1}, y_{1}) \in \mathrm{E}_{1}^{2}\)。由于 \(\varrho_{2}\) 是酉的且 q 不变，得到

\[
b \left(\varrho_ {1} (g) x _ {1}, \varrho_ {1} (g) y _ {1}\right) = q \left(x _ {1} \otimes \varrho_ {2} \left(g ^ {- 1}\right) \varepsilon , y _ {1} \otimes \varrho_ {2} \left(g ^ {- 1}\right) \varepsilon\right) = b \left(x _ {1}, y _ {1}\right),
\]

因此 \(b \in \mathrm{Sesq}_{\mathrm{G}_1}(\varrho_1, \varrho_1)\)。于是存在唯一的 \(\lambda \in \mathbf{C}\)，使得 \(b(x_1, y_1) = \lambda \langle x_1 | y_1 \rangle\) 对所有 \((x_1, y_1) \in \mathrm{E}_1^2\) 成立（推论 3），即

\[
q (x _ {1} \otimes x _ {2}, y _ {1} \otimes y _ {2}) = \lambda \langle x _ {1} | y _ {1} \rangle \langle x _ {2} | y _ {2} \rangle
\]

对所有 \((x_{1},y_{1},x_{2},y_{2})\in\mathrm{E}_{1}^{2}\times\mathrm{E}_{2}^{2}\) 成立。由此可推出空间 \(\operatorname{Sesq}_{\mathrm{G}_{1}\times\mathrm{G}_{2}}(\varrho,\varrho)\) 的维数为 1，这意味着表示 \(\varrho=\varrho_{1}\boxtimes\varrho_{2}\) 不可约（同上）。

回忆一下，U 表示模为 1 的复数所成的群。

推论 7. --- 设 \(\pi\) 是拓扑群 G 在希尔伯特空间 E 中的不可约酉表示。存在从 G 的中心 C 到 U 的连续同态 \(\chi\)，满足 \(\pi(z) = \chi(z) \cdot 1_{\mathrm{E}}\) 对所有 \(z \in \mathrm{C}\) 成立。特别地，若 G 是交换群，则 \(\dim(\mathrm{E}) = 1\)。

对于 \(z \in \mathrm{C}\)，映射 \(\pi(z)\) 属于 \(\mathrm{Hom}_{\mathrm{G}}(\pi, \pi)\)；因此它具有 \(\chi(z) \cdot 1_{\mathrm{E}}\) 的形式，其中 \(\chi(z)\) 是某个复数（命题 6）。由于 \(\pi(z)\) 是酉映射，有 \(|\chi(z)| = 1\)。此外，由于 \(\pi\) 是同态，映射 \(z \mapsto \chi(z)\) 是同态。取 E 中 \(v \neq 0\)；从 C 到 U 的映射 \(z \mapsto \chi(z)v = \pi(z)v\) 是连续的，因此同态 \(\chi\) 是连续的。

最后，若 G 是交换群，则 C = G，所以对 G 中所有 g，有 \(\pi(g) = \chi(g) \cdot 1_{\mathrm{E}}\)；于是 E 的每个一维子空间都是 \(\pi\) 的子表示，而由于 \(\pi\) 不可约，空间 E 必为一维。

定义 6. --- 设 π 是拓扑群 G 在希尔伯特空间 E 中的不可约酉表示。满足对 C 中所有 z 有 π(z) = χ(z) · \(1_{\mathrm{E}}\) 的从 G 的中心 C 到 U 的同态 χ，称为 π 的中心特征标。

注. --- 设 \(\pi\)（分别为 \(\varrho\)）是拓扑群 G（分别为 H）在希尔伯特空间 E（分别为 F）中的不可约酉表示。记 \(\chi\)（分别为 \(\eta\)）为 \(\pi\)（分别为 \(\varrho\)）的中心特征标。表示 \(\overline{\pi}\) 的中心特征标是 \(\overline{\chi}\)；而 G × H 的不可约酉表示 \(\pi \boxtimes \varrho\)（推论 6）的中心特征标，是 G × H 上的特征标 \(\chi \boxtimes \eta\)：\((g, h) \mapsto \chi(g)\eta(h)\)。

\subsection{半单性}

定义 7. --- 设 \(\varrho\) 是拓扑群 G 在希尔伯特空间 E 中的酉表示；称 \(\varrho\) 为半单的，当且仅当存在不可约子表示族 \((\mathrm{F}_i)_{i\in \mathrm{I}}\)，它们是 \(\varrho\) 的子表示，且 E 是子空间 \(\mathrm{F}_i\) 的希尔伯特直和。

有时也说半单酉表示具有离散分解，或称其可离散分解。

如果 \((\varrho_{i})_{i\in\mathrm{I}}\) 是拓扑群 G 的半单酉表示族，那么表示 \(\varrho_{i}\) 的希尔伯特直和是半单的。

命题 7. --- 设 G 是拓扑群，\(\varrho\) 是 G 在复希尔伯特空间 E 中的酉表示。表示 \(\varrho\) 半单，当且仅当 \(\varrho\) 的每个非零子表示都包含不可约子表示。

假设 \(\varrho\) 半单。设 \((\mathrm{F}_i)_{i\in \mathrm{I}}\) 是 \(\varrho\) 的不可约子表示族，使得 E 是子空间 \(\mathrm{F}_i\) 的希尔伯特直和。设 F 是 E 的一个非零子表示。存在 \(i\in \mathrm{I}\)，使得像为 \(\mathrm{F}_i\) 的正交投影在 F 上的限制非零。该正交投影于是通过限制到子空间定义了非零 G-态射 \(p_i\)，从 F 到 \(\mathrm{F}_i\)。由舒尔引理，该 G-态射 \(p_i\) 是满射（V, p. 388 的推论 5）。F 中 \(p_i\) 的核的正交补是 F 的一个与 \(\mathrm{F}_i\) 同构的子表示（V, p. 383 的命题 3），因而不可约。

证明逆命题。令 \(\mathcal{F}\) 为 E 的闭子空间 F 所成的集合，其中 F 在 G 下稳定，且 \(\varrho\) 在 F 中诱导的子表示不可约。令 \(\mathcal{O}\) 为所有满足 \(\mathcal{G}\) 是 \(\mathcal{F}\) 的子集且其成员两两正交的 \(\mathcal{G}\) 所成的集合。

集合 \(\mathcal{O}\) 具有有限特征（E, III, p. 34, 定义 2）。于是令 \(\mathcal{G}\) 为 \(\mathcal{O}\) 的一个极大元（E, III, p. 35, 定理 1）。令 E \(_1\) 为 E 中的子空间，它是 \(\mathcal{G}\) 中不可约子表示 F 的希尔伯特直和。若 E \(_1\) ≠ E，则 E \(_1\) 的正交补非零，而根据假设，G 在 E \(_1^{\circ}\) 中的子表示包含一个不可约子表示。其表示空间 F 与 \(\mathcal{G}\) 中的元素正交，于是 \(\mathcal{G} \cup \{F\} \in \mathcal{O}\)；这与 \(\mathcal{G}\) 的极大性矛盾。因此 E = E \(_1\)，证明完毕。

推论 1. --- 设 G 是拓扑群，\(\varrho\) 是 G 在复希尔伯特空间 E 中的半单酉表示。\(\varrho\) 的每个子表示（以及 \(\varrho\) 的每个商表示）都是半单的。

如果 \(\varrho_{1}\) 是 \(\varrho\) 的子表示，那么 \(\varrho_{1}\) 的每个非零子表示都包含不可约子表示（命题 7），故 \(\varrho_{1}\) 是半单的（同上）。

由 V, p. 383 的命题 3，\(\varrho\) 的每个商表示都与 \(\varrho\) 的一个子表示同构；因此它是半单的。

推论 2. --- 拓扑群 G 的每个有限维酉表示 \(\varrho\) 都是半单的。

这由 V, p. 379 的命题 7 和引理 3 得出。

推论 3. --- 设 \(\varrho_{1}\) 和 \(\varrho_{2}\) 是拓扑群 G 的有限维酉表示。表示 \(\varrho_{1}\) 与 \(\varrho_{2}\) 同构，当且仅当 \(\chi_{\varrho_{1}} = \chi_{\varrho_{2}}\)。

这由推论 2 以及 A, VIII, p. 389 的命题 1, b) 得出。

\subsection{酉表示的类}

引理 9. --- 设 E 是 K 上的希尔伯特空间，F ⊂ E 是 E 的稠密向量子空间。那么 E 的希尔伯特维数小于或等于 F 的维数。

若 E 的希尔伯特维数有限，则它等于 E 的维数。设 E 的希尔伯特维数无限，则子空间 F 的维数也是无限的。设 B 是 E 的标准正交基，B' 是 F 的基。对于每个 \(x \in B\)，存在 \(f(x) \in B'\)，使得 \(\langle x \mid f(x) \rangle \neq 0\)，否则将有 \(x \in F^{\circ} = \{0\}\)。于是定义映射 \(f: B \to B'\)。

对于每个 \(y \in B'\)，集合 \(f^{-1}(y)\) 包含于满足 \(x \in B\) 且 \(\langle x | y \rangle \neq 0\) 的集合中，因而是可数的（EVT, V, p. 21, 命题 4）。由 E, III, p. 50, 命题 4，得到 \(\text{Card(B)} = \text{Card}(f(B)) \leqslant \text{Card(B')}\)。

记 \(Is_{G}(\pi_{1},\pi_{2})\) 为关系

“G 是拓扑群，且 \(\pi_{1}\) 和 \(\pi_{2}\) 是

K 上希尔伯特空间中的酉表示，并且二者

同构。”关于 \(\pi_{1}\) 和 \(\pi_{2}\)，这是一个等价关系。对于 K 上希尔伯特空间中 G 的任意酉表示 \(\pi\)，记 \(\mathrm{cl}(\pi)\) 为 \(\pi\) 的等价类（参见 E, II, p. 47）；因此它是 G 的一个与 \(\pi\) 同构的酉表示；称 \(\mathrm{cl}(\pi)\) 为 \(\pi\) 的类。K 上希尔伯特空间中的酉表示 \(\pi_{1}\) 和 \(\pi_{2}\) 同构，当且仅当 \(\mathrm{cl}(\pi_{1}) = \mathrm{cl}(\pi_{2})\)。

设 G 是拓扑群，c 是一个基数。关系

“λ 是 G 在

希尔伯特维数 \(\leqslant\mathfrak{c}\) 的复希尔伯特空间中的酉表示的类”

是关于 \(\lambda\) 的集化关系（E, II, p. 3）。事实上，每个维数 \(\leqslant \mathfrak{c}\) 的 K 上希尔伯特空间都等距同构于 \(\ell^2 (\mathfrak{c})\) 的一个希尔伯特子空间（EVT, V, p. 23, cor. 2），于是断言由 E, II, p. 47 得出。

设 \(\pi\) 是 G 在希尔伯特空间 E 中的不可约酉表示。取 E 中的非零元素 \(x\)。由于 \(\pi\) 不可约，向量 \(x\) 是 \(\pi\) 的循环向量，这意味着 E 的希尔伯特维数 \(\leqslant\) Card(G)（将引理 9 应用于由 \(\pi(g)x\)（\(g \in G\)）所生成的稠密子空间）。因此，G 在 K 上希尔伯特空间中的不可约酉表示的各类，属于 G 在 K 上希尔伯特维数 \(\leqslant\) Card(G) 的希尔伯特空间中的酉表示的类所成的集合；所以它们构成一个集合。

定义 8. --- 我们用 \(\widehat{\mathbf{G}}\) 表示 \(\mathbf{G}\) 在复希尔伯特空间中的不可约酉表示的类所成的集合。称 \(\widehat{\mathbf{G}}\) 为 \(\mathbf{G}\) 的酉对偶。

因此，对于 G 在复希尔伯特空间中的每个不可约酉表示 \(\pi\)，都有 \(\operatorname{cl}(\pi) \in \widehat{\mathbf{G}}\)。

备注. --- 1) 假设 G 是交换群。由 V, p. 390 的推论 7，G 在复希尔伯特空间中的每个不可约酉表示都是一维的。如果 G 局部紧，则集合 \(\widehat{G}\) 可与 G 的酉特征标集合等同（II, p. 201 的定义 1），所以记号 \(\widehat{G}\) 与 II, p. 201 的定义 2 中引入的记号相容。

\begin{enumerate}
\def\labelenumi{\arabic{enumi})}
\setcounter{enumi}{1}
\item
  如果 G 是有限群，则集合 \(\widehat{\mathbf{G}}\) 与 \(\mathbf{C}[\mathrm{G}]\)-单模的类所成的集合之间存在双射（A, VIII, p. 47），后者在 A, VIII, p. 396 中也记作 \(\widehat{\mathbf{G}}\)。
\item
  对于 \(\pi \in \widehat{\mathbf{G}}\)，我们把 \(\overline{\pi}\) 与 \(\widehat{\mathbf{G}}\) 中 \(\pi\) 的共轭表示的类等同。
\item
  如果 \(\pi\) 是 G 在有限维复希尔伯特空间中的不可约酉表示，则其特征标 \(\chi_{\pi}\) 只取决于 \(\pi\) 的类。因此可以谈论 G 的有限维复不可约酉表示的特征标集合。
\end{enumerate}

\subsection{同型分量}

定义 9. --- 设 G 是拓扑群，\(\pi\) 是 G 的连续不可约表示。设 \(\varrho\) 是 G 在分离局部凸空间 E 中的连续表示。称 \(\pi\)-同型分量（对应 \(\varrho\)）为 \(\varrho\) 的所有与 \(\pi\) 同构的子表示的表示空间之和在 E 中的闭包。该子空间记为 \(\mathrm{M}_{\pi}(\varrho)\)。

空间 \(\mathrm{M}_{\pi}(\varrho)\) 是 E 的闭子空间，并定义 \(\varrho\) 的一个子表示。该空间只取决于 \(\pi\) 在 \(\widehat{G}\) 中的类。

命题 8. — 设 G 是拓扑群，\(\varrho\) 是 G 在复希尔伯特空间 E 中的酉表示。

\begin{enumerate}
\def\labelenumi{\alph{enumi})}
\item
  设 \(\pi_1\) 和 \(\pi_2\) 是 G 的非同构不可约酉表示。子空间 \(\mathrm{M}_{\pi_1}(\varrho)\) 和 \(\mathrm{M}_{\pi_2}(\varrho)\) 正交；
\item
  设 \(\varrho'\) 是 G 的酉表示，且 \(u\) 是从 \(\varrho\) 到 \(\varrho'\) 的 G-态射。对于 G 的每个不可约酉表示 \(\pi\)，有 \(u(\mathrm{M}_{\pi}(\varrho)) \subset \mathrm{M}_{\pi}(\varrho')\)。
\end{enumerate}

设 \(E_{1}\) 和 \(E_{2}\) 是 E 的子空间，分别定义与 \(\pi_{1}\) 和 \(\pi_{2}\) 同构的子表示。E 上像为 \(E_{2}\) 的正交投影通过限制到子空间定义了 \(\mathrm{Hom}_{\mathrm{G}}(\pi_{1},\pi_{2})\) 的一个元素；该元素为零（V, p. 387 的推论 2），这证明了 \(E_{2}\) 与 \(E_{1}\) 正交。断言 \(a\) 得证。

对于断言 \(b\)，注意到根据定义，\(\mathrm{M}_{\pi}(\varrho)\) 是 E 中由如下元素生成的空间的闭包：元素 \(x \in \mathrm{E}\) 属于闭子空间 \(\mathrm{F} \subset \mathrm{E}\)，该子空间在 \(\varrho\) 下稳定，且其子表示 \(\varrho_{\mathrm{F}}\) 是 \(\varrho\) 的子表示并与 \(\pi\) 同构。因此只需证明在这种情形下 \(u(x) \in \mathrm{M}_{\pi}(\varrho')\)。令 \(\mathrm{H} = u(\mathrm{F})\)；它是 \(\varrho'\) 的表示空间的闭子空间（V, p. 388 的推论 5），在 \(\varrho'\) 下稳定，而由限制到子空间得到的 u 是从 F 到 H 的满射 G-态射。如果 H 非零，则由 V, p. 387 的推论 2，该 G-态射是同构。因此 G 在 H 上的表示与 \(\pi\) 同构，从而 \(u(x)\) 属于 \(\mathrm{M}_{\pi}(\varrho')\)。

对于任意向量空间 H 及 H 的子空间族 \((\mathrm{H}_{i})_{i\in\mathrm{I}}\)，称这些空间 \((\mathrm{H}_{i})_{i\in\mathrm{I}}\) 为直和，如果族 \((\mathrm{H}_{i})_{i\in\mathrm{I}}\) 满足 A, II, p. 18 的命题 11 的等价条件。

命题 9. --- 设 G 是拓扑群，\(\varrho\) 是 G 在复希尔伯特空间 E 中的酉表示。设 \(\pi\) 是 G 在复希尔伯特空间 \(\mathrm{E}_{\pi}\) 中的不可约酉表示。记 \(v\) 为从 \(\mathrm{Hom}_{\mathrm{G}}(\pi, \varrho) \otimes \mathrm{E}_{\pi}\) 到 E 的线性映射，满足 \(v(u \otimes x) = u(x)\) 对所有 \((u, x) \in \mathrm{Hom}_{\mathrm{G}}(\pi, \varrho) \times \mathrm{E}_{\pi}\) 成立。

线性映射 \(v\) 是单射，其像是 \(\varrho\) 的所有与 \(\pi\) 同构的子表示的表示空间之和。特别地，\(v\) 的像在 \(M_{\pi}(\varrho)\) 中稠密。

V, p. 388 的推论 5 表明，v 的像是 \(\varrho\) 的所有与 \(\pi\) 同构的子表示的表示空间之和。

证明 \(v\) 是单射。设 \((u_i)_{i \in I}\) 是向量空间 Hom \(_G(\pi, \varrho)\) 的一组基。对于 \(i \in I\)，记 F \(_i\) 为 \(u_i\) 的像，记 \(\widetilde{u}_i \colon E_\pi \to F_i\) 为由 \(u_i\) 通过限制到子空间得到的 G-同构。

首先对 I 的有限子集 J 的基数作归纳，证明子空间 \((\mathrm{F}_{i})_{i\in\mathrm{J}}\) 是直和。

J 为空的情形是显然的。设 J 非空，并且所需性质对 I 中基数至多为 \(\text{Card(J)} - 1\) 的子集成立。

取 J 中的一个固定元素 \(j\)。归纳假设说明子空间 \(\mathrm{F}_{i}\)（\(i\in\mathrm{J}-\{j\}\)）构成直和。令 \(\mathrm{F}^{\prime}\) 为它们的和；现在只需证明 \(\mathrm{F}_{j}\cap\mathrm{F}^{\prime}=\{0\}\)。

假设 \(F_{j}\) 与 \(F'\) 的交不为 0；该交是 \(F_{j}\) 的一个子表示，而由于后者不可约，故有 \(F_{j} \cap F' = F_{j}\)，从而 \(F_{j} \subset F'\)。

对于 \(i \in \mathrm{J} - \{j\}\)，记 \(\operatorname{pr}_i: \mathrm{F}' \to \mathrm{F}_i\) 为投影，记 \(\iota_i: \mathrm{F}_i \to \mathrm{F}'\) 为包含映射。我们有规范同构

\[
\operatorname{Hom} _ {\mathrm{G}} \left(\mathrm{F} _ {j}, \mathrm{F} ^ {\prime}\right)\rightarrow \bigoplus_ {i \in \mathrm{J} - \{j \}} \operatorname{Hom} _ {\mathrm{G}} \left(\mathrm{F} _ {j}, \mathrm{F} _ {i}\right)
\]

（公式（1）, p. 377）使得 \(u \colon \mathrm{F}_j \to \mathrm{F}'\) 的像为族 \((\mathrm{pr}_i \circ u)_{i \in \mathrm{J} - \{j\}}\)（A, II, p. 13, cor. 1）。V, p. 387 的推论 2 表明，非零的 \(\widetilde{u}_i \circ \widetilde{u}_j^{-1}\) 是空间 \(\mathrm{Hom}_{\mathrm{G}}(\mathrm{F}_j, \mathrm{F}_i)\) 的一组基。

因此，G-态射族 \((\iota_{i} \circ \widetilde{u}_{i} \circ \widetilde{u}_{j}^{-1})_{i \in \mathrm{J} - \{j\}}\) 是 \(\mathrm{Hom}_{\mathrm{G}}(\mathrm{F}_{j}, \mathrm{F}')\) 的一组基。

记 \(\iota\) 为 \(F_{j}\) 到 \(F'\) 的包含映射；它是 \(\mathrm{Hom}_{\mathrm{G}}(\mathrm{F}_{j}, \mathrm{F}')\) 的一个元素，因而是各 G-态射 \(\iota_{i} \circ \widetilde{u}_{i} \circ \widetilde{u}_{j}^{-1}\)（\(i \in \mathrm{J} - \{j\}\)）的线性组合。于是 \(u_{j}\) 是各映射 \(u_{i}\)（\(i \neq j\)）的线性组合，这与族 \((u_{i})_{i \in \mathrm{I}}\) 的线性无关性矛盾。因此断言由归纳得到证明。

现在证明 v 是单射。设 w 是 \(\mathrm{Hom}_{\mathrm{G}}(\pi,\varrho)\otimes\mathrm{E}_{\pi}\) 的一个元素。存在唯一的有限支集族 \((x_{i})_{i\in\mathrm{I}}\)，其元素属于 \(E_{\pi}\)，使得

\[
w = \sum_ {i \in \mathrm{I}} u _ {i} \otimes x _ {i}
\]

（A, II, p. 62, cor. 1），于是有

\[
v (w) = \sum_ {i \in I} u _ {i} (x _ {i}).
\]

由前述内容，条件 \(v(w) = 0\) 因而蕴含 \(u_{i}(x_{i}) = 0\) 对所有 \(i \in I\) 成立，因而对所有 i 有 \(x_{i} = 0\)，这是因为 \(u_{i}\) 是单射，从而 w = 0。

命题 10. --- 设 G 是拓扑群。设 \(\pi\) 是 G 在复希尔伯特空间 E \(_{\pi}\) 中的不可约酉表示，\(\varrho\) 是 G 在复希尔伯特空间 E 中的酉表示。

存在 E 的闭不变子空间族 \((\mathrm{E}_{i})_{i\in\mathrm{I}}\)，使得 \(\varrho\) 在 \(E_{i}\) 中的子表示与 \(\pi\) 同构，对所有 \(i\in I\) 成立，并且 \(\mathrm{M}_{\pi}(\varrho)\) 是子空间 \(E_{i}\) 的希尔伯特直和。此外，满足这些性质的族 \((\mathrm{E}_{i})_{i\in\mathrm{I}}\) 的指标集 I 的基数与所选族无关。

证明满足所述条件的族存在。令 \(\mathcal{O}\) 为 \(\mathrm{Hom}_{\mathrm{G}}(\pi, \varrho)\setminus\{0\}\) 的子集 C 所成的集合，使得子空间 \(u(\mathrm{E}_{\pi})\)（\(u \in \mathrm{C}\)）两两正交。集合 \(\mathcal{O}\) 按包含关系排序。它非空，因为空集属于其中；它具有有限特征（E, III, p. 34, 定义 2），因为 C 属于 \(\mathcal{O}\) 当且仅当 C 的每个二元子集属于 \(\mathcal{O}\)。由 E, III, p. 35, 定理 1，\(\mathcal{O}\) 存在极大元 C。令 F 为由 \(u(\mathrm{E}_{\pi})\)（\(u \in \mathrm{C}\)）所成子空间生成的子空间的闭包；它是 \(u(\mathrm{E}_{\pi})\)（\(u \in C\)）所成空间的希尔伯特直和。我们将证明 \(\mathrm{F} = \mathrm{M}_{\pi}(\varrho)\)，从而证明族 \((u(\mathrm{E}_{\pi}))_{u \in \mathrm{C}}\) 具有所需性质。

按定义，\(u(\mathrm{E}_{\pi}) \subset \mathrm{M}_{\pi}(\varrho)\) 对所有 \(u \in \mathrm{C}\) 成立，故 \(\mathrm{F}\) 包含于 \(\mathrm{M}_{\pi}(\varrho)\)。为证明反向包含，只需证明：若 \(v\) 是从 \(\pi\) 到 \(\varrho\) 的 G-态射，则其像包含于 \(\mathrm{F}\)。令 \(p\) 为 \(\mathrm{E}\) 上像为 \(\mathrm{F}^{\circ}\) 的正交投影；由于 \(\mathrm{F}\) 是 \(\mathrm{E}\) 的子表示，它是 G-态射（V, p. 383 的命题 4）。G-态射 \(p \circ v\) 的像与 \(\mathrm{F}\) 正交；因此该像为零（否则 \(\mathrm{C} \cup \{p \circ v\} \in \mathcal{O}\)，这与 \(\mathrm{C}\) 的极大性矛盾），从而 \(v\) 的像包含于 \(\mathrm{F}\)。

现在设 \((\mathrm{E}_{i})_{i\in\mathrm{I}}\) 和 \((\mathrm{F}_{j})_{j\in\mathrm{J}}\) 是 E 的两两正交闭不变子空间族，使得 G 在 \(E_{i}\)（分别在 \(F_{j}\)）中的子表示与 \(\pi\) 同构，对所有 \(i\in I\)（分别对每个 \(j\in J\)）成立，并且 \(\mathrm{M}_{\pi}(\varrho)\) 既是族 \((\mathrm{E}_{i})_{i\in\mathrm{I}}\) 的希尔伯特直和，也是族 \((\mathrm{F}_{j})_{j\in\mathrm{J}}\) 的希尔伯特直和。

若 I 有限，则

\[
\dim \operatorname{Hom} _ {\mathrm{G}} (\pi , \mathrm{M} _ {\pi} (\varrho)) = \dim \operatorname{Hom} _ {\mathrm{G}} \Bigl (\mathrm{E} _ {\pi}, \bigoplus_ {i \in \mathrm{I}} \mathrm{E} _ {i} \Bigr) = \operatorname{Card} (\mathrm{I})
\]

（V, p. 377 的公式（1）和 V, p. 387 的推论 2）。于是对于 J 的每个有限子集 L，有

\[
\begin{array}{c} \operatorname{Card} (\mathrm{L}) = \dim \operatorname{Hom} _ {\mathrm{G}} \Bigl (\mathrm{E} _ {\pi}, \bigoplus_ {j \in \mathrm{L}} \mathrm{F} _ {j} \Bigr) \\ \leqslant \dim \operatorname{Hom} _ {\mathrm{G}} (\pi , \operatorname{M} _ {\pi} (\varrho)) = \operatorname{Card} (\mathrm{I}) \end{array}
\]

（同上）。这证明了 J 也是有限的，且 \(\text{Card(I)} = \text{Card(J)}\)，正如所要求的。同样，若 J 有限，则 I 有限且具有相同的基数。

现在假设 I 和 J 都是无限的。对于 \(j \in J\)，记 \(p_{j}\) 为 E 上像为 \(F_{j}\) 的正交投影。对于 \(i \in I\)，取 \(x_{i}\) 为 \(E_{i}\) 中的非零元素；它是 \(E_{i}\) 的循环向量。注意，由于 \(p_{j}\) 通过限制到子空间得到的是从 \(E_{i}\) 到 \(F_{j}\) 的 G-态射，空间 \(p_{j}(E_{i})\) 为零，当且仅当 \(p_{j}(x_{i}) = 0\)（事实上，空间 \(E_{i} \cap \text{Ker}(p_{j})\) 要么为零，要么等于 \(E_{i}\)）。

对于每个 \(j \in J\)，存在 \(f(j) \in I\)，使得 \(p_{j}(\mathrm{E}_{f(j)})\) 不为 0（否则，正交投影 \(p_{j}\) 在 \(\mathrm{M}_{\pi}(\varrho)\) 上为零）。于是定义了从 J 到 I 的映射 f。对于每个 \(i \in I\)，集合 \(f^{-1}(i)\) 是可数的。事实上，该集合包含于由 \(j \in J\) 中满足 \(p_{j}(\mathrm{E}_{i})\) 非零者组成的集合中，也就是满足 \(p_{j}(x_{i}) \neq 0\) 的那些 j。现在由（EVT, V, p. 20, cor. 2），有

\[
\sum_ {j \in \mathrm{J}} \| p _ {j} (x _ {i}) \| ^ {2} = \| x _ {i} \| ^ {2},
\]

因此，\(j \in J\) 中满足 \(p_{j}(x_{i}) \neq 0\) 的元素所成集合是可数的。由 E, III, p. 50, 命题 4，得到 \(\text{Card}(\text{J}) = \text{Card}(f(\text{J})) \leqslant \text{Card}(\text{I})\)。交换 I 和 J 的角色，得到 \(\text{Card}(\text{I}) = \text{Card}(\text{J})\)。

推论. --- 设 G 是拓扑群，\(\varrho\) 是 G 在复希尔伯特空间 E 中的酉表示。G 的 \(\pi\)-同型分量（其中 \(\pi \in \widehat{\mathbf{G}}\)）的希尔伯特直和是 \(\varrho\) 的最大半单子表示。

事实上，\(\pi\)-同型分量对每个 \(\pi \in \widehat{\mathbf{G}}\) 两两正交（命题 8，\(a\)），且每个都是半单的（命题 10，\(a\)），所以空间 \(\mathrm{M}_{\pi}(\varrho)\) 对 \(\pi \in \widehat{\mathbf{G}}\) 的希尔伯特直和 F 定义了 \(\varrho\) 的一个半单子表示。此外，\(\varrho\) 的每个不可约子表示都是 \(\varrho\) 的某个同型分量的子表示，故推论得证。

定义 10. --- 设 G 是拓扑群。设 \(\varrho\) 是 G 在复希尔伯特空间 E 中的酉表示，\(\pi\) 是 G 在复希尔伯特空间中的不可约酉表示。

称 \(\pi\) 在 \(\varrho\) 中的重数为指标集的基数，所针对的是 E 的在 G 下稳定的闭子空间族 \((\mathrm{E}_i)_{i\in \mathrm{I}}\)，其中 \(\varrho\) 在 \(\mathrm{E}_i\) 中诱导的子表示与 \(\pi\) 同构，对每个 \(i\in I\) 成立，并且 \(\mathrm{M}_{\pi}(\varrho)\) 是子空间 \(\mathrm{E}_i\) 的希尔伯特直和。

如果 \(\pi\) 在 \(\varrho\) 中的重数有限，则由 V, p. 377 的公式（1）和 V, p. 387 的推论 2，它等于空间 \(\mathrm{Hom}_{\mathrm{G}}(\pi, \varrho)\) 的维数（分别等于 \(\mathrm{Hom}_{\mathrm{G}}(\varrho, \pi)\) 的维数）。一般而言这并不总是成立。

注. --- 一个酉表示 \(\varrho\) 可能非零，但相对于 G 的所有不可约表示的 \(\varrho\) 的所有同型分量都为零（参见 V, p. 426, 注）。

\section{局部紧群的表示}

本段中，向量空间均在域 K = R 或 C 上。记 G 为赋予左哈尔测度 \(\mu\) 的局部紧群。记 \(\mathcal{L}^{p}(\mathrm{G}) = \mathcal{L}_{\mathbf{C}}^{p}(\mathrm{G}, \mu)\)，\(\mathrm{L}^{p}(\mathrm{G}) = \mathrm{L}_{\mathbf{C}}^{p}(\mathrm{G}, \mu)\)，其中 \(p \in [1, +\infty]\)。

\subsection{某些表示的连续性}

命题 1. — 设 H 是局部紧群。设 \(\varrho_{1}\)（分别为 \(\varrho_{2}\)）是 G（分别为 H）在分离局部凸 K-向量空间 E \(_{1}\)（分别为 E \(_{2}\)）中的连续表示。令 F 表示赋予紧收敛拓扑的空间 \(\mathcal{L}(\mathrm{E}_{1};\mathrm{E}_{2})\)。由 \(\varrho\) 通过 \(\varrho(g,h)u=\varrho_{2}(h)\circ u\circ\varrho_{1}(g^{-1})\)（\((g,h)\in\mathrm{G}\times\mathrm{H}\)）定义的 G×H 到 F 的表示是连续的。

记 \(\mathcal{L}_c(\mathrm{E}_1;\mathrm{E}_2)\) 为赋予紧收敛拓扑的空间 \(\mathcal{L}(\mathrm{E}_1;\mathrm{E}_2)\)。设 \(\mathfrak{F}_1\)（分别为 \(\mathfrak{F}_2,\mathfrak{F}_3\)）是趋于 \(e\) 的 G 中滤子（分别为趋于 0 的 \(\mathcal{L}_c(\mathrm{E}_1;\mathrm{E}_2)\) 中滤子、趋于 \(e\) 的 H 中滤子）。由于 G 和 H 局部紧，存在 \(\mathrm{C}\in \mathfrak{F}_1\) 和 \(\mathrm{D}\in \mathfrak{F}_3\)，它们是相对紧的。集合 \(\varrho_1(\mathrm{C}^{-1})\) 在 \(\mathcal{L}(\mathrm{E}_1)\) 中等度连续（参见 INT, VIII, p. 129, § 2, no. 1, 注 2, \(a'\)）。由 EVT, III, p. 33 的命题 9 和 EVT, III, p. 31 的命题 4，由 \((u,v)\mapsto u\circ v\) 定义的从 \(\mathcal{L}_c(\mathrm{E}_1;\mathrm{E}_2)\times \varrho_1(\mathrm{C}^{-1})\) 到 \(\mathcal{L}_c(\mathrm{E}_1;\mathrm{E}_2)\) 的映射是连续的。因此，滤基 \(\mathfrak{F}_2\circ \varrho_1(\mathfrak{F}_1^{-1})\) 在 \(\mathcal{L}_c(\mathrm{E}_1;\mathrm{E}_2)\) 中趋于 0。集合 \(\varrho_2(\mathrm{D})\) 在 \(\mathcal{L}(\mathrm{E}_2)\) 中等度连续（参见 INT, VIII, p. 129, § 2, n° 1, rem. 2），并且对每个 \(x\in \mathrm{E}_2\)，集合 \(\varrho_2(\mathrm{D})x\subset \mathrm{E}_2\) 是相对紧的。因此，赋予紧收敛拓扑的 \(\varrho_2(\mathrm{D})\) 在 \(\mathcal{L}(\mathrm{E}_2)\) 中相对紧（TG, X, p. 18, cor. 1）。由 \((u,v)\mapsto u\circ v\) 定义的从 \(\overline{\varrho_2(\mathrm{D})}\times \mathcal{L}_c(\mathrm{E}_1;\mathrm{E}_2)\) 到 \(\mathcal{L}_c(\mathrm{E}_1;\mathrm{E}_2)\) 的映射，由 EVT, III, p. 33 的命题 9 和 EVT, III, p. 31 的命题 4，是连续的。因此，滤基 \(\varrho (\mathfrak{F}_1\times \mathfrak{F}_3)(\mathfrak{F}_2)\) 在 \(\mathcal{L}_c(\mathrm{E}_1;\mathrm{E}_2)\) 中趋于 0。这就证明了断言。

推论. --- 设 \(\varrho\) 是 G 在分离局部凸 K-向量空间 E 上的连续表示。当 E' 赋予紧收敛拓扑时，逆步表示 \(\breve{\varrho}\) 是连续的。

取 \(H = \{e\}\)、\(\varrho_{1} = \varrho\)，并令 \(\varrho_{2}\) 为 K 上的平凡表示，即由命题得到。

注. --- 当 E' 赋予强拓扑时，逆步表示不一定连续（参见 INT, VIII, p. 191, § 2, 习题 3, d)）。可以证明，如果空间 E 是半自反的，则它是连续的（同上，c)）。

\subsection{表示向测度空间的延拓}

本节中假设 K = C。

设 \(\varrho\) 是 G 在分离局部凸拟完备空间 E 中的连续线性表示。记 \(\mathcal{M}_{c}(\mathrm{G})\) 为 G 上具有紧支撑且赋予紧收敛拓扑的测度空间；它是空间 \(\mathcal{C}(\mathrm{G})\) 的对偶。对于每个测度 \(\nu \in \mathcal{M}_{c}(\mathrm{G})\)，置

\[
\varrho (\nu) = \int_ {\mathrm{G}} \varrho (g) d \nu (g).
\]

这是 \(\mathcal{L}(\mathrm{E})\) 的一个元素。特别地，\(\varrho(\varepsilon_g) = \varrho(g)\) 对所有 \(g \in \mathrm{G}\) 成立。

根据 INT, VIII, p. 136, § 2, no. 6，映射 \(\nu \mapsto \varrho(\nu)\) 是从 \(\mathcal{M}_c(\mathrm{G})\) 到赋予紧收敛拓扑的空间 \(\mathcal{L}(\mathrm{E})\) 的连续线性映射。根据 INT, VIII, p. 145, § 3, no. 3, prop. 11，它是含单位代数的态射。

设 \(x \in \mathrm{E}\)。当 \(\mathcal{M}_c(\mathrm{G})\) 赋予紧收敛拓扑时，由 \(\nu \mapsto \varrho(\nu)x\) 定义的从 \(\mathcal{M}_c(\mathrm{G})\) 到 E 的映射是连续的（INT, VI, p. 27, § 1, no. 7, prop. 16 以及 EVT, III, p. 31, prop. 4）。

对于在 G 的某个紧子集之外为零的 \(f \in \mathcal{L}^{1}(\mathrm{G})\)，测度 \(f \cdot \mu\) 具有紧支撑；我们把 E 的自同态 \(\varrho^{\mu}(f)\)（也记为 \(\varrho(f)\)）记为 \(\varrho(f \cdot \mu)\)。此自同态只取决于其类 \(\widetilde{f}\)，即 \(f\) 在 \(\mathrm{L}^{1}(\mathrm{G})\) 中的类，我们也把它记为 \(\varrho^{\mu}(\widetilde{f})\) 或 \(\varrho(\widetilde{f})\)。

同样，设 \(\varrho\) 是 G 在巴拿赫空间 E 中的连续有界线性表示。对于每个有界测度 \(\nu \in \mathcal{M}^{1}(\mathrm{G})\)，置

\[
\varrho (\nu) = \int_ {\mathrm{G}} \varrho (g) d \nu (g).
\]

根据 INT, VIII, 同上，映射 \(\nu \mapsto \varrho(\nu)\) 是从 G 上有界复测度代数 \(\mathcal{M}^{1}(G)\) 到巴拿赫代数 \(\mathcal{L}(E)\) 的连续含单位巴拿赫代数态射。

令 \(\rho\) 为 G 上的函数 \(g\mapsto\|\varrho(g)\|\)；INT, VIII, p. 145 的命题 11 中记作 \(\mathcal{M}^{\rho}\) 的代数（其元素是测度 \(\nu\)，满足 \(\rho\in\mathcal{L}^{1}(\nu)\)）与巴拿赫代数 \(\mathcal{M}^{1}(G)\) 重合。事实上，置 \(M = \sup \rho\)。有 \(M > 0\)，因为 \(\rho(e)=1\)，且 \(M^{-1} \leqslant \rho \leqslant M\)，因为不等式 \(\|\varrho(e)\|\leqslant\|\varrho(g)\|\|\varrho(g^{-1})\|\) 对所有 \(g\in G\) 成立。因此，\(\rho\in\mathcal{L}^{1}(\nu)\) 当且仅当 \(\nu\) 是有界测度。

若 \(f \in \mathcal{L}^{1}(\mathrm{G})\)，自同态 \(\varrho^{\mu}(f)\)（或简记为 \(\varrho(f)\)）就是 \(\varrho(f \cdot \mu)\)；对于类 \(\tilde{f}\)，它是 \(f\) 在 \(\mathrm{L}^{1}(\mathrm{G})\) 中的类，同样如此。

引理 1. --- 设 \(\varrho\) 是 G 在希尔伯特空间 E 中的酉表示。映射 \(\nu \mapsto \varrho(\nu)\) 从 \(\mathcal{M}^{1}(\mathrm{G})\) 到 \(\mathcal{L}(\mathrm{E})\)，是对合巴拿赫代数的含单位态射。

由前述内容，只需证明态射 \(\nu \mapsto \varrho(\nu)\) 是对合的。设 \(\nu \in \mathcal{M}^{1}(\mathrm{G})\)。测度 \(\nu^{*}\) 是测度 \(\check{\nu}\) 的共轭测度（I, p. 99, 例 4）。根据共轭测度的定义（INT, III, p. 52, § 1, no. 5），计算得

\[
\begin{array}{c} \langle x \mid \varrho (\nu) y \rangle = \int_ {\mathrm{G}} \langle x \mid \varrho (g) y \rangle d \nu (g) \\ = \int_ {\mathrm{G}} \langle \varrho (g ^ {- 1}) x \mid y \rangle d \nu (g) = \int_ {\mathrm{G}} \langle \varrho (g) x \mid y \rangle d \check {\nu} (g) \\ = \overline {{\int_ {\mathrm{G}} \langle y \mid \varrho (g) x \rangle d \nu^ {*} (g)}} = \langle \varrho (\nu^ {*}) x \mid y \rangle \end{array}
\]

对 E 中所有 \(x\) 和 \(y\) 成立，从而 \(\varrho(\nu)^{*} = \varrho(\nu^{*})\)。

设 \(\varrho\) 是 G 在局部凸拟完备空间 E（分别在巴拿赫空间 E）中的连续线性表示（分别为连续有界表示）。如果 F 是 E 的闭子空间，并且子表示 \(\varrho_{\mathrm{F}}\) 由 \(\varrho\) 在 F 上诱导，那么对于每个测度 \(\nu \in \mathcal{M}_c(\mathrm{G})\)（分别为 \(\nu \in \mathcal{M}^1 (\mathrm{G})\)），线性映射 \(\varrho (\nu)\) 保持子空间 F 不变，并且在 F 上与 \(\varrho_{\mathrm{F}}(\nu)\) 重合。

反之，设 \(F \subset E\) 是闭子空间，并且在每个线性映射 \(\varrho(f)\) 下稳定，其中 \(f \in \mathcal{H}(G)\)。则 F 定义 \(\varrho\) 的一个子表示（INT, VIII, p. 139, § 2, no. 7, prop. 10）。

回忆（A, VIII, p. 388），函数 \(f \colon \mathbf{G} \to \mathbf{C}\) 称为中心函数，是指对所有 \(g\) 和 \(h\)（均属于 \(\mathbf{G}\)），有 \(f(gh) = f(hg)\)。这等价于说 \(f\) 在共轭作用下不变。

定义 1. --- 有界测度 \(\nu \in \mathcal{M}^{1}(\mathrm{G})\) 称为中心的，是指 \(\varepsilon_{g} * \nu = \nu * \varepsilon_{g}\) 对所有 \(g \in \mathrm{G}\) 成立。

如果 G 是单模群，且 \(f \in \mathcal{C}(\mathrm{G})\) 关于 \(\mu\) 可积，那么测度 \(f \cdot \mu\) 中心，当且仅当 f 是中心函数。

设 \(\varrho\) 是 G 在巴拿赫空间 E 中的连续有界表示（分别是 G 在拟完备分离局部凸空间 E 中的连续表示）。对于 G 上的每个有界中心测度 \(\nu\)（分别为 G 上的任意紧支撑中心测度 \(\nu\)），线性映射 \(\varrho(\nu)\) 是从 \(\varrho\) 到 \(\varrho\) 的 G-态射。事实上，对每个 \(g \in G\)，有

\[
\varrho (\nu) \varrho (g) = \varrho (\nu * \varepsilon_ {g}) = \varrho (\varepsilon_ {g} * \nu) = \varrho (g) \varrho (\nu).
\]

\subsection{半单性判据}

命题 2. --- 设 \(\varrho\) 是 G 在复希尔伯特空间 E 中的酉表示。设 \(\mathfrak{F}\) 是 \(\mathcal{M}_{c}(\mathrm{G})\) 上的滤基，并且对于紧收敛拓扑趋于测度 \(\varepsilon_{e}\)。假设存在 \(\mathrm{M} \in \mathfrak{F}\)，使得 \(\varrho(\nu)\) 对于每个 \(\nu \in \mathrm{M}\) 都是 E 的紧自同态。

那么酉表示 \(\varrho\) 是半单的，并且 G 的每个不可约酉表示在 \(\varrho\) 中的重数有限。

先证明一个引理。

引理 2. --- 设 \(\varrho_{1}\) 是 G 的一个非零酉表示，并且与 \(\varrho\) 的一个子表示同构。存在测度 \(\nu \in M\)，使得 \(\varrho_{1}(\check{\nu} * \nu)\) 是 \(\varrho_{1}\) 的表示空间上的非零紧埃尔米特自同态。特别地，存在非零实数 \(\lambda\)，使得 \(\varrho_{1}(\check{\nu} * \nu)\) 关于 \(\lambda\) 的特征空间非零。

不妨设 \(\varrho_{1}\) 是 \(\varrho\) 的子表示。令 E \(_{1}\) 为其表示空间。对于每个测度 \(\nu \in \mathrm{M}\)，自同态 \(\varrho_{1}(\nu)\) 由 \(\varrho(\nu)\) 限制到子空间得到。因此，根据假设及 III, p. 5 的命题 3，\(\varrho_{1}(\nu)\) 是紧的。

由于 \(\mathfrak{F}\) 趋于测度 \(\varepsilon_{e}\)，且收敛所在的空间是赋予紧收敛拓扑的 \(\mathcal{M}_{c}(\mathrm{G})\)，而空间 \(\mathrm{E}_{1}\) 非零，存在 \(\nu\in\mathrm{M}\)，使得 \(v=\varrho_{1}(\nu)\) 是 \(\mathrm{E}_{1}\) 的非零自同态（参见 V, p. 400, n°2）。自同态 \(u=\varrho_{1}(\check{\nu}*\nu)\) 等于 \(v^{*}\circ v\)（V, p. 401 的引理 1）；因此它非零（EVT, V, p. 39, prop. 2），是埃尔米特的且是紧的，因为 \(v\) 是紧的。

由于 \(u\) 是非零埃尔米特自同态，其谱不只含 0（I, p. 110 的例 1）。最后，由于 \(u\) 是紧的，每个 \(\lambda \in \mathrm{Sp}(u)\setminus\{0\}\) 都是 \(u\) 的特征值（III, p. 83 的命题 2）。

现在证明命题。

我们将使用 V, p. 391 的命题 7 来证明 \(\varrho\) 半单。设 \(\varrho_{1}\) 是 \(\varrho\) 的一个非零子表示，\(\mathrm{E}_{1} \subset \mathrm{E}\) 是其表示空间；我们必须证明 \(\varrho_{1}\) 包含一个不可约子表示。

取 \(\nu \in \mathrm{M}\)，使得 \(u = \varrho_{1}(\check{\nu}*\nu)\) 是 \(\varrho_{1}\) 的表示空间上的非零紧埃尔米特自同态，并取非零的 \(\lambda\)，使得 \(u\) 关于 \(\lambda\) 的特征空间 F 非零（引理 2）。空间 F 是有限维的（III, p. 90 的命题 5）。存在子表示 \(\varrho_{2}\)，它是 \(\varrho_{1}\) 在子空间 \(\mathrm{E}_{2}\)（\(\mathrm{E}_{1}\) 的子空间）上的表示，使得 \(\mathrm{E}_{2} \cap \mathrm{F}\) 非零且维数最小。取非零元素 \(x\)，它属于 \(\mathrm{E}_{2} \cap \mathrm{F}\)；令 \(\mathrm{E}_{3}\) 为 \(\varrho_{1}\) 中由 \(x\) 生成的子表示。于是 \(\mathrm{E}_{3} \subset \mathrm{E}_{2}\)，故 \(\mathrm{E}_{3} \cap \mathrm{F} = \mathrm{E}_{2} \cap \mathrm{F}\)，这是由 \(\mathrm{E}_{2} \cap \mathrm{F}\) 的维数极小性得到的。

证明表示 \(E_{3}\) 不可约。设 \(E_{4} \subset E_{3}\) 是在 G 下稳定的闭子空间。有 \(E_{2} \cap F = (E_{4} \cap F) \oplus (E_{4}^{\circ} \cap F)\)，其中 \(E_{4}^{\circ}\) 表示 \(E_{4}\) 在空间 \(E_{3}\) 中的正交补（事实上，若 \(y \in E_{2} \cap F\)，令 \(y_{4} \in E_{4}\) 为 y 在 \(E_{4}\) 上的正交投影；由于 \(E_{4}\) 和 \(E_{4}^{\circ}\) 在 u 下稳定，向量 \(u(y_{4})\) 是 \(u(y) = \lambda y\) 在 \(E_{4}\) 上的正交投影，从而 \(u(y_{4}) = \lambda y_{4}\)，这意味着 \(y_{4} \in E_{4} \cap F\)，断言得证）。于是由 \(E_{2} \cap F\) 的维数极小性，或者有 \(E_{4} \cap F = E_{2} \cap F\)，或者有 \(E_{4}^{\circ} \cap F = E_{2} \cap F\)。在第一种情形中，\(x \in E_{4}\)，从而 \(E_{4} = E_{3}\)；在第二种情形中，\(x \in E_{4}^{\circ}\)，从而 \(E_{4}^{\circ} = E_{3}\) 且 \(E_{4} = \{0\}\)。因此表示 \(E_{3}\) 不可约。

由 V, p. 391 的命题 7，得到表示 \(\varrho\) 是半单的。

设 \(\pi\) 是 G 的不可约酉表示，且其在 \(\varrho\) 中的重数非零；令 \(E_{\pi}\) 为其表示空间。存在测度 \(\nu \in M\) 和非零实数 \(\lambda\)，使得 \(u = \pi (\check{\nu} * \nu)\) 是 \(E_{\pi}\) 上的非零紧埃尔米特自同态，并且 \(u\) 关于 \(\lambda\) 的特征空间非零（引理 2）。令 F 为 \(v = \varrho (\check{\nu} * \nu)\) 关于 \(\lambda\) 的特征空间。对于每个子表示 \(E_1\)，它是 E 的与 \(\pi\) 同构的子表示；在 \(E_1\) 上由 \(v\) 限制到子空间得到的自同态与 \(u\) 等同，且 \(E_1\) 与 F 的交非零。因此，\(\pi\) 在 \(\varrho\) 中的重数小于空间 F 的维数，而该维数有限（III, p. 90 的命题 5）。

例. --- 假设 G 是单模群，例如实半单李群。设 \(\Gamma \subset G\) 是离散子群，使得商空间 \(X = \Gamma \backslash G\) 紧。群 G 通过右乘作用在 X 上。记 \(\beta\) 为 \(\Gamma\) 上的计数测度，并置 \(\widetilde{\mu} = \mu / \beta\)（INT, VII, p. 44, § 2, n° 2, 定义 1）；这是 X 上的有界 G-不变测度。

对于每个 \(f \in \mathcal{L}^2(\mathrm{X},\widetilde{\mu})\) 和每个 \(g \in \mathrm{G}\)，定义函数 \(\varrho(g)f \in \mathcal{L}^2(\mathrm{X},\widetilde{\mu})\)，使其满足 \(\varrho(g)f(x) = f(x \cdot g)\)。映射 \(\varrho(g)\) 是连续线性映射，通过取商在 \(\mathrm{L}^2(\mathrm{X},\widetilde{\mu})\) 上诱导出一个酉映射，仍记为 \(\varrho(g)\)。映射 \(\varrho\) 是 G 在 \(\mathrm{L}^2(\mathrm{X},\widetilde{\mu})\) 中的酉表示（INT, VII, p. 135, § 2, n° 5, prop. 8）。

取 \(\varphi \in \mathcal{K}(\mathrm{G})\)。对于 \(\mathrm{L}_1\) 和 \(\mathrm{L}_2\) 这两个 \(\mathrm{G}\) 的任意紧子集，\(\mathrm{T}\) 是 \(\Gamma\) 与紧集 \(\mathrm{L}_1\mathrm{Supp}(\varphi)\mathrm{L}_2^{-1}\) 的交，且该交是有限的。对于每个 \((g,h) \in \mathrm{L}_1 \times \mathrm{L}_2\)，级数 \(\sum_{\gamma \in \Gamma} \varphi(g^{-1}\gamma h)\) 与有限和 \(\sum_{\gamma \in \mathrm{T}} \varphi(g^{-1}\gamma h)\) 重合。由于 \(\mathrm{G}\) 局部紧，该级数的和 \(k_{\varphi}(g,h)\) 是 \(\mathrm{G} \times \mathrm{G}\) 上的连续函数。

取 \((g,h) \in G \times G\) 以及 \((\gamma, \eta) \in \Gamma \times \Gamma\)。有 \(k_{\varphi}(\gamma g, \eta h) = k_{\varphi}(g, h)\)，因此 \(k_{\varphi}\) 通过取商定义了 \(X \times X\) 上的连续函数，记为 \(\widetilde{k}_{\varphi}\)。又有 \(\widetilde{k}_{\varphi} \in \mathcal{L}^{2}(X \times X, \widetilde{\mu} \otimes \widetilde{\mu})\)，因为假设空间 X 是紧的。

记 \(\dot{x}\) 为 G 中元素 \(x\) 在规范投影下于 X 中的像。取 \(\varphi \in \mathcal{K}(\mathrm{G})\)。对于 \(f \in \mathcal{L}^2(\mathrm{X},\widetilde{\mu})\) 和 \(x \in \mathrm{G}\)，有

\[
\begin{array}{c} \varrho (\varphi) f (\dot {x}) = \int_ {\mathrm{G}} \varphi (g)   (\varrho (g) f) (\dot {x}) d \mu (g) = \int_ {\mathrm{G}} \varphi (g) f (\dot {x} \cdot g) d \mu (g) \\ = \int_ {\mathrm{G}} \varphi (x ^ {- 1} y) f (\dot {y}) d \mu (y) = \int_ {\Gamma \backslash \mathrm{G}} \Bigl (\sum_ {\gamma \in \Gamma} \varphi (x ^ {- 1} \gamma y) \Bigr) f (\dot {y}) d \widetilde {\mu} (\dot {y}) \end{array}
\]

（INT, VII, p. 46, § 2, no. 3, prop. 5）。由于 \(\widetilde{k}_{\varphi}\) 属于空间 \(\mathcal{L}^2 (\mathrm{X}\times \mathrm{X},\widetilde{\mu}\otimes \widetilde{\mu})\)，所以自同态 \(\varrho (\varphi)\) 在 \(\mathrm{L}^2 (\mathrm{X},\widetilde{\mu})\) 上与核为 \(\widetilde{k}_{\varphi}\) 的希尔伯特–施密特自同态重合；因此该自同态是紧的（III, p. 33 的推论 1）。

存在函数序列 \((\varphi_{n})_{n\in\mathbf{N}}\)，其元素属于 \(\mathcal{K}_{+}(\mathrm{G})\) 且积分为 1，使得 \(\varphi_{n}\cdot\mu\) 趋于测度 \(\varepsilon_{e}\)（在赋予紧收敛拓扑的 \(\mathcal{M}_{c}(\mathrm{G})\) 中）（INT, VIII, p. 139, § 2, no. 7, cor. 2）。因此命题 2 蕴含表示 \(\varrho\) 半单，且 G 的不可约酉表示在 \(\varrho\) 中的重数有限。

在 \(\varrho\) 中重数非零的 G 的不可约酉表示，称为群 G 的 \(\Gamma\)-自守表示。

\subsection{正则表示}

设 \(p\) 是满足 \(\geqslant 1\) 的实数，\(\mu'\) 是 \(G\) 上的右哈尔测度。对于每个 \(g \in G\) 及每个函数 \(f \in \mathcal{L}^p(G, \mu)\)，记 \(\gamma_G^{(p)}(g)f\) 为函数 \(x \mapsto f(g^{-1}x)\)（定义在 \(G\) 上）。映射 \(g \mapsto \gamma_G^{(p)}(g)\) 是 \(G\) 在 \(\mathcal{L}^p(G)\) 中的连续线性表示。通过取商，它诱导出 \(G\) 在 \(L^p(G)\) 中的连续等距线性表示（INT, VIII, p. 135, § 2, no. 5, prop. 8），记为 \(\gamma_G^{(p)}\)。

同样，对于每个 \(g \in G\) 及每个函数 \(f \in \mathcal{L}^{p}(G, \mu')\)，记 \(\boldsymbol{\delta}_{\mathrm{G}}^{(p)}(g)f\) 为 G 上的函数 \(x \mapsto f(xg)\)。映射 \(g \mapsto \boldsymbol{\delta}_{\mathrm{G}}^{(p)}(g)\) 是 G 在 \(\mathcal{L}^{p}(G, \mu')\) 中的连续线性表示。通过取商，它诱导出 G 在 \(\mathrm{L}^{p}(\mathrm{G}, \mu')\) 中的连续等距线性表示（参见 INT, VIII, p. 136, § 2, no. 5）。

称 \(\boldsymbol{\gamma}_{\mathrm{G}}^{(p)}\) 为 G 在 \(\mathcal{L}^{p}(\mathrm{G})\) 或 \(\mathrm{L}^{p}(\mathrm{G})\) 中的左正则表示，称 \(\boldsymbol{\delta}_{\mathrm{G}}^{(p)}\) 为 G 在 \(\mathcal{L}^{p}(\mathrm{G},\mu^{\prime})\) 或 \(\mathrm{L}^{p}(\mathrm{G},\mu^{\prime})\) 中的右正则表示。

引理 3. — 设 p 是大于等于 1 的实数。G 在 L \(^{p}\) (G,μ) 中的左（分别为右）正则表示（分别在 L \(^{p}\) (G,μ') 中）是忠实的。

更确切地，令 q 为 p 的共轭指数，令 g 为 G 中满足 \(g \neq e\) 的元素。

\begin{enumerate}
\def\labelenumi{\alph{enumi})}
\item
  存在函数 \(\varphi\in\mathcal{K}(\mathrm{G})\)，它在 \(\mathrm{L}^{q}(\mathrm{G},\mu)\) 中为正且非零，并满足 \(\langle\varphi,\boldsymbol{\gamma}_{\mathrm{G}}^{(p)}(g)\varphi\rangle=0\)；
\item
  存在函数 \(\varphi\in\mathcal{K}(\mathrm{G})\)，它在 \(\mathrm{L}^{q}(\mathrm{G},\mu')\) 中为正且非零，并满足 \(\langle\varphi,\boldsymbol{\delta}_{\mathrm{G}}^{(p)}(g)\varphi\rangle=0\)。
\end{enumerate}

考虑左正则表示 \(\boldsymbol{\gamma}_{\mathrm{G}}^{(p)}\) 的情形，右正则表示的情形类似。断言 \(a\) 蕴含 \(\boldsymbol{\gamma}_{\mathrm{G}}^{(p)}\) 是忠实的，因为若 \(g \neq e\) 是 G 中的元素，且 \(\varphi\) 如 \(a\) 所述，则有 \(\varphi \neq \boldsymbol{\gamma}_{\mathrm{G}}^{(p)}(g)\varphi\)，因为 \(\langle \varphi, \varphi \rangle = \int_{\mathrm{G}} \varphi^2 > 0\)。

因此证明 \(a\)。令 \(g \neq e\) 属于 G。令 C 是 \(e\) 的一个紧对称邻域，使得 \(g \notin \mathrm{C}^2\)，并令 \(\varphi \in \mathcal{K}(\mathrm{G})\) 是积分为 1 且支撑包含于 C 的连续正函数；函数 \(\varphi\) 在 \(\mathrm{L}^p(\mathrm{G}, \mu)\) 中非零。由于 \(\mathrm{C} \cap g^{-1}\mathrm{C} = \emptyset\)，有

\[
\langle \varphi , \boldsymbol {\gamma} _ {\mathrm{G}} ^ {(p)} (g) \varphi \rangle = \int_ {\mathrm{G}} \varphi (x) \varphi (g ^ {- 1} x) d \mu (x) = 0.
\]

设 \(\varrho\) 是 G 在巴拿赫空间 E 中的连续有界表示。对于每个 \(f\in \mathrm{L}^1 (\mathrm{G})\)（分别为 \(f^{\prime}\in \mathrm{L}^{1}(\mathrm{G},\mu^{\prime}))\)）以及每个 \(g\in \mathrm{G}\)，有

\[
\varrho (g) \varrho (f \cdot \mu) = \varrho (\varepsilon_ {g} * (f \cdot \mu)) = \varrho (\boldsymbol {\gamma} _ {\mathrm{G}} ^ {(1)} (g) f \cdot \mu),\tag{1}
\]

\[
\varrho (f ^ {\prime} \cdot \mu^ {\prime}) \varrho (g) = \varrho ((f ^ {\prime} \cdot \mu^ {\prime}) * \varepsilon_ {g}) = \varrho (\boldsymbol {\delta} _ {\mathrm{G}} ^ {(1)} (g ^ {- 1}) f ^ {\prime} \cdot \mu ^ {\prime}),\tag{2}
\]

（INT, VIII, p. 144, § 3, no. 2, 公式（5））。

若 G 是单模群，置 \(\mu' = \mu\)。G 在 \(\mathcal{L}^{p}(\mathrm{G})\)（分别在 \(\mathrm{L}^{p}(\mathrm{G})\)）中的双正则表示 \(\varrho_{\mathrm{G}}^{(p)}\)，是 \(G \times G\) 在 \(\mathcal{L}^{p}(\mathrm{G})\)（分别在 \(\mathrm{L}^{p}(\mathrm{G})\)）中的表示，满足

\[
\boldsymbol {\varrho} _ {\mathrm{G}} ^ {(p)} (g _ {1}, g _ {2}) = \boldsymbol {\gamma} _ {\mathrm{G}} ^ {(p)} (g _ {1}) \boldsymbol {\delta} _ {\mathrm{G}} ^ {(p)} (g _ {2}) = \boldsymbol {\delta} _ {\mathrm{G}} ^ {(p)} (g _ {2}) \boldsymbol {\gamma} _ {\mathrm{G}} ^ {(p)} (g _ {1}).
\]

这是连续线性表示（V, p. 377 的引理 1）。它满足

\[
(\boldsymbol {\varrho} _ {\mathrm{G}} ^ {(p)} (g _ {1}, g _ {2}) f) (x) = f (g _ {1} ^ {- 1} x g _ {2})
\]

对所有 \(f \in \mathcal{L}^p(\mathrm{G})\)、所有 \((g_1, g_2) \in \mathrm{G} \times \mathrm{G}\) 及每个 \(x \in \mathrm{G}\) 成立。

注. --- G 在 \(\mathrm{L}^{p}(\mathrm{G},\mu)\) 中的双正则表示不一定忠实；其核是 G 的中心在映射 \(g\mapsto(g,g)\) 下的像（V, p. 487 的习题 4）。

当 \(p = 2\) 时，G 的左正则表示 \(\gamma_{\mathrm{G}}^{(2)}\) 在复希尔伯特空间 \(\mathrm{L}^2 (\mathrm{G},\mu)\) 中是酉的，因为它是等距的。同样，右正则表示 \(\delta_{\mathrm{G}}^{(2)}\) 在 \(\mathrm{L}^2 (\mathrm{G},\mu')\) 中也是酉的。

我们简记 \(\gamma_{\mathrm{G}} = \gamma_{\mathrm{G}}^{(2)}\) 和 \(\delta_{\mathrm{G}} = \delta_{\mathrm{G}}^{(2)}\)，并将这些表示称为 G 的左、右正则表示。

若 G 是单模群，则双正则表示 \(\varrho_{\mathrm{G}}^{(2)}\) 是 \(G \times G\) 在 \(\mathrm{L}^{2}(\mathrm{G}, \mu)\) 中的表示，且是酉的。我们简记它为 \(\varrho_{G}\)。

引理 4. --- 设 \(p\) 是满足 \(\geqslant 1\) 的实数。对于 \(f \in \mathcal{L}^1(\mathrm{G})\)，线性映射 \(\gamma_{\mathrm{G}}^{(p)}(f)\) 与 \(\mathrm{L}^p(\mathrm{G})\) 上由 \(\varphi \mapsto f *^\mu \varphi\) 定义的自同态重合。

这由 INT, VIII, p. 157, § 4, n°2 的命题 6 得出，同时要注意 INT, VIII, p. 165 的公式（14）。

注. --- 回忆一下，对于 G 上的任意函数 \(f\)，在 G 上定义函数 \(\check{f}\)，使得 \(\check{f}(g) = f(g^{-1})\) 对所有 \(g \in \mathrm{G}\) 成立（INT, VII, p. 12, § 1, n° 1, 公式（12））。

可以检验，若 \(f \in \mathcal{L}^{1}(\mathrm{G})\)，线性映射 \(u = \boldsymbol{\delta}_{\mathrm{G}}^{(p)}(f)\) 满足关系式

\[
\widetilde {u (\varphi)} = f * \check {\varphi}
\]

对所有 \(\varphi\in\mathrm{L}^{p}(\mathrm{G},\mu^{\prime})\) 成立。